%% notes-tex/019-simb-multimodal/main.tex
%% [[experiments.019-simb-multimodal.phenotype-strand-retrospective]]
%% https://github.com/Mjvolk3/torchcell/tree/main/notes-tex/019-simb-multimodal/main.tex
%%
%% Typeset counterpart of the dendron hierarchy notes/experiments.019-simb-multimodal.*.md,
%% plus the 012 / 020 / 021 / 022 / 023 experiment notes it draws on. Pairing is by NAME:
%% notes-tex/<slug>/ <-> notes/<slug>.*.md.
%%
%% Build:  make            -> main.pdf         (draft: status chips + provenance flags)
%%         make clean-view -> main-clean.pdf   (share view)
%%         make check      -> formatting / provenance gate
%%         make plots-all  -> refresh figures/ from the plot scripts' true-size SVGs
%%         make figures    -> refresh figures/ from draw.io sources
%%
%% EVERY number in this document comes from a committed script under experiments/ or
%% from a named W&B run id, and is marked with a `%% SOURCE:` comment.

\documentclass[10pt,a4paper]{article}
\usepackage[draft]{tcdoc}

\title{\vspace{-1.5em}Phenotype Strands After SIMB\\
  {\large Expression, Morphology, and Metabolite Production:
   What Was Measured and What to Run Next}}
\author{Michael Volk}
\date{\today}

\begin{document}
\maketitle
\thispagestyle{empty}

\begin{abstract}
\noindent
Six phenotype strands were run against one shared cell embedding: knockout expression,
morphology, the two jointly, betaxanthin production, free amino-acid pools, and
beta-carotene production. Two were presented at SIMB 2026. Here is what each measured,
against what ceiling, under what scoring rule, read from all 28 W\&B projects and 2,187
runs rather than from a selection.

Every strand sits further from its ceiling than the manuscript's placeholder text assumes,
and the reason differs by strand. \textbf{Expression} is under-trained rather than
saturated: no run has peaked at any budget up to 10,000 epochs, and the score tracks the
epoch budget across a factor of 500 in budget. Its training loss and its metric come apart,
and the culprit is specifically the \emph{quantile} loss, which bottoms at epoch 463 and
rises for the next 9,500 while Pearson climbs; squared error does not diverge, bottoming
alongside the Pearson peak. Separately, the model orders genes at $r \approx 0.24$ while
sitting level with ``predict each gene's mean'' in squared error, because its predictions are
$1.8\times$ more spread out than that correlation justifies; one free rescale takes it
clearly past the mean predictor without changing any correlation.
\textbf{Morphology} has never been trained on more than a quarter of its screen, in any of
397 runs across eight projects, and its ceiling is measured here for the first time at
$0.611$. \textbf{Betaxanthin} is the strongest strand, at $47\,\%$ of a $0.914$ ceiling, and
its argument is coverage rather than accuracy: a flux-based method has features for
$19\,\%$ of the screen and misses $73\,\%$ of its high producers. The \textbf{free
amino-acid} profile of a deletion predicts that deletion's betaxanthin at out-of-fold
$r = 0.298$ while tyrosine, the chromophore precursor, predicts it at $0.064$, so a real
coupling exists on an axis that is not the precursor, and a plain auxiliary head does not
use it. \textbf{Beta-carotene} is limited by its measurement, a hand-scored ordinal with one
replicate per strain. And the \textbf{graph prior}, reported at chance in the previous
revision, is not: asked the question that is defined for every phenotype rather than for
expression alone, adjacency reaches AUC $0.66$ against a $0.50$ control, and which graph
carries signal differs by phenotype.

The campaign that follows is expression first, single-phenotype and at full budget, because
it is the strand where the next experiment is cheapest to interpret and the diagnosis is
furthest along. Morphology at full scale is deferred rather than dropped.
\end{abstract}

\vspace{0.5em}
{\small\color{tcgray}
\textbf{How to read this.} \cref{sec:terms} defines the words that recur. \cref{sec:summary}
is the whole argument and every number a decision turns on.
\crefrange{sec:expression}{sec:carotene} are one section per strand, each ending in what it
established and what it cannot say. \cref{sec:common} collects the failures the strands
share, which is where the leverage is. \cref{sec:directions} ranks the options, and
\cref{sec:campaign} sizes the next campaign. \textbf{Going straight to the campaign:}
\cref{sec:objective}, \cref{tab:morph-ntrain} and \cref{sec:campaign}.\par\smallskip

Revision 3 answers 127 review comments on revision 2, comment by comment. The ledgers are\\
\file{notes-tex/019-simb-multimodal/review/round-2-dispositions.md} (this round) and\\
\file{notes-tex/019-simb-multimodal/review/round-1-dispositions.md} (the previous one).
Claims that revision 2 got wrong are corrected in place and named in the ledger, not
argued away.\par\smallskip

Working notes: \file{notes/experiments.019-simb-multimodal.*.md},
\file{notes/experiments.020-cachera-betaxanthin.*.md}\par
Generating scripts: the \file{scripts/} directory of\\
\file{experiments/019-simb-multimodal/}, \file{experiments/020-cachera-betaxanthin/}, and\\
\file{experiments/023-metabolome-betaxanthin-joint/}\par
Leaderboard: \file{experiments/019-simb-multimodal/results/round_leaderboards.csv}\par}

\vspace{1em}
\tccontents

\clearpage
%% notes-tex/019-simb-multimodal/sections/1-summary.tex
%% SOURCE: every number here is carried from a later section and cited there. Nothing
%% originates in this section.

\section{Summary\secstatus{todo}}
\label{sec:summary}

\subsection{Terms used throughout\secstatus{todo}}
\label{sec:terms}

\begin{description}
\item[Head] The output module for one phenotype. The encoder and the perturbation operator
  are shared; a head is the last few layers that turn the shared representation into one
  phenotype's values. A run declares which heads are active, so \emph{betaxanthin head} and
  \emph{morphology head} name output modules, not models.
\item[Metabolome head] The head that predicts the 19 Mulleder amino-acid concentrations.
  It appears in two roles. As the \textbf{only} active head it is the amino-acid strand
  (\cref{sec:amino-acid}). Added \textbf{alongside} the betaxanthin head it is an
  \term{auxiliary head}: it is trained and scored, but the number being reported is still
  betaxanthin, and the question is whether sharing an encoder with it helps.
\item[Auxiliary head] Any head added to a run for its effect on a different head's score.
  Adding one changes two things at once, the gradient signal and the population of strains
  the run can use, which is why every comparison here pins the instance set.
\item[Strand] One phenotype prosecuted across a series of W\&B projects.
\item[\file{pearson_per_feature}] For each output feature (a reporter gene, a CalMorph
  parameter, one amino acid), the correlation across strains between prediction and
  measurement; then averaged over features. It goes to zero when a model predicts each
  feature's mean. The companion \file{pearson_per_instance} correlates across features
  within one strain and stays high under exactly that collapse, so it is never the
  objective here.
\item[\file{nmse}] Mean squared error divided by the variance of the target about each
  feature's mean, so $\texttt{nmse} = 1$ is exactly ``predict each gene's training mean''.
\item[\file{roll_max}] The scoring rule: the maximum of a centered five-point rolling mean
  of a validation curve. An upward-biased order statistic whose bias grows with the number
  of epochs run, so every score here travels with its epoch count.
\item[Ceiling] The highest correlation any predictor could reach against a target carrying
  measurement noise, estimated per strand from replicate structure.
\end{description}

\subsection{Where each strand stands\secstatus{todo}}

%% GENERATED by experiments/019-simb-multimodal/scripts/build_retrospective_tables.py
%% SOURCE: results/round_leaderboards.csv + the three ceiling JSONs. Do not hand-edit.
\begin{table}[htbp]
\centering\small
\begin{tabular}{lrrrrl}
\toprule
strand & ceiling & best & of ceiling & epochs (peak/last) & runs \\
\midrule
Expression (Kemmeren, Sameith) & 0.775 & \textbf{0.229} & 30\% & 4007/4105 & 163 \\
Expression, masked-label objective & 0.775 & \textbf{0.241} & 31\% & 9188/9997 & 16 \\
Morphology (CalMorph, Ohya) & 0.611 & \textbf{0.082} & 13\% & 27/65 & 23 (1 flat) \\
Expression and morphology, joint & -- & \textbf{0.051} & -- & 28/70 & 32 \\
Betaxanthin (Cachera) & 0.914 & \textbf{0.372} & 41\% & 254/999 & 44 \\
Betaxanthin with metabolome head & 0.914 & \textbf{0.430} & 47\% & 176/496 & 46 \\
Amino-acid pools (Mulleder) & -- & \textbf{0.211} & -- & 402/999 & 31 \\
Beta-carotene (Ozaydin) & 0.544 & \textbf{0.132} & 24\% & 799/999 & 39 \\
\bottomrule
\end{tabular}
\caption[]{Every strand, scored the same way. \emph{best} is \texttt{roll\_max}, the maximum of a centered five-point rolling mean of the validation curve, an upward-biased order statistic whose bias grows with epochs run, which is why the epoch column travels with it. \emph{epochs (peak/last)} gives where the best run peaked and where it stopped: a peak in the last tenth means the run was cut off while still improving. Ceilings are measured by three different estimators because the three measurement types require them; the amino-acid cell is empty because Mulleder has one replicate per strain, and no ceiling is estimable. \emph{flat} counts runs whose validation metric never left zero.}
\label{tab:strand-summary}
\end{table}


Read the epoch column with the score column. Expression and beta-carotene peak in the last
tenth of their best run, so those numbers are lower bounds set by the compute budget.
Morphology and the expression-morphology joint arm peak at epochs 27 and 28 of runs that
stopped at 65 and 70, which is not a result of any kind. Betaxanthin and amino acid peak in
the first half of a 999-epoch run and then overfit.

The two betaxanthin rows are not comparable to each other: the joint row trains on 3,832
strains against the plain row's 3,698, because the joint configuration restricts to
genotypes carrying both phenotypes. The within-grid paired comparison in
\cref{tab:bx-aa-paired} is the valid form of that contrast, and it comes out negative.

\subsection{Nine findings\secstatus{todo}}

\begin{enumerate}
\item \textbf{The expression strand contains no valid comparison between mechanisms.} All
  67 scored arms stopped at or below 300 epochs, against a validation curve that dips at
  200 to 300 and then climbs past its early peak. No peak has been observed at any budget
  up to 10,000 epochs. Every per-arm delta on record measures early-training transients.

\item \textbf{The expression objective is fighting itself, and it is the quantile loss
  rather than squared error.} On the longest run the quantile loss bottoms at epoch 463 and
  rises for the next 9,500 while Pearson climbs, but squared error bottoms at epoch 9,175,
  alongside the Pearson peak at 9,674; on the second-longest run the two are one epoch
  apart. \file{nmse} never returns below 1 after about epoch 400, so the model reaches
  $r = 0.236$ while being \emph{no better than predicting each gene's mean} in squared
  error. The cause is calibration, not capability: at the peak its predictions are
  $1.95\times$ more spread out than its own correlation justifies, and a post-hoc rescale
  by $r/s$ moves \file{nmse} from $1.010$ to $0.944$ while changing no correlation at all.
  Best-by-\file{mse} checkpointing would have picked within 500 epochs of the right model,
  so checkpointing is a narrower problem than previously recorded and is not what holds the
  strand at $0.24$.

\item \textbf{One expression run is a multi-day object, and this sets the campaign.} The
  best run took \textbf{91.4 hours, 3.81 days}, for 9,999 epochs at 32.9 s/epoch, and still
  peaked at $97\,\%$ of the way through. One GPU delivers $\approx 2{,}500$ epochs per day,
  so the number of arms a campaign can afford is fixed before any of them are chosen.

\item \textbf{The perturbation operator provably cannot express a pair term, and masked
  unmasking does not supply one.} At one perturbed gene the cross-attention softmax runs
  over a one-element key set, so the model reduces to $g(h_i + c_b)$ and nothing depends on
  which reporter is predicted given which gene was deleted. 16,200 of 32,760 attention
  parameters are dead. The scGPT-style scheme conditions on \emph{other genes' measured
  values}, and at $k = 0$, which is where scoring happens, every one of those features is
  exactly zero and the forward pass is identical to the unconditioned model. The two are
  different axes. Capacity is not the binding constraint either: a rank-32 reconstruction
  of the target reaches $0.727$. A rank-64 pair term already reaches $0.2274$ in 2.0 days
  against the masked objective's $0.2362$ in 3.8.

\item \textbf{Morphology has never been trained on its own data.} The Ohya screen holds
  4,718 deletions; every morphology run used the 1,440 that also carry expression, giving
  $n_{\mathrm{train}} = 1{,}161$. Its ceiling, measured here for the first time, is
  $\mathbf{0.611}$ over 278 features, with 201 of them individually above $0.5$. For a
  second-label experiment, fitness shares 4,220 strains with morphology against
  expression's 1,440.

\item \textbf{The free amino-acid profile predicts betaxanthin; the precursor does not.}
  Over the 4,432 deletions the two screens share, the 19-dimensional pool predicts
  betaxanthin at out-of-fold $r = \mathbf{0.298}$, while tyrosine alone, the chromophore
  precursor, predicts it at $\mathbf{0.064}$ and correlates marginally at $-0.076$. Roughly
  three quarters of the profile's power survives regressing out single-mutant fitness. On
  the 724 deletions with no Costanzo fitness record, which carry twice the betaxanthin
  spread, the profile predicts at $0.455$.

\item \textbf{The metabolome auxiliary head currently costs betaxanthin performance.}
  Paired within grid cell, joint minus control is $-0.0265 \pm 0.0159$ over the five
  comparable cells where the control learned the task, with 2 of 5 positive. The
  $+0.0203$ that motivated the replication reproduces, at $+0.0192$, in the one cell it came
  from. So the coupling exists in the data and the current head does not use it.

\item \textbf{The graph prior is at chance, and the mask deletes the one signal that
  exists.} Graph proximity does not predict which reporters respond to a deletion, on any of
  the nine graphs, in the orientation the model uses: AUC $0.4961$ to $0.5057$, at most
  $+0.0046$ over a degree-preserving control, and longer walks make it worse. Three of the
  nine graphs are also silent on most deletions, scoring only $29$ to $75\,\%$ of them. But
  \emph{direction} carries real signal: following tflink edges TF $\to$ target gives
  $0.5508$ and regulatory interaction $0.5239$, against controls near $0.501$. The mask
  builder symmetrizes, which averages that away. Nine attention heads are currently spent
  enforcing a target that does not predict response.

\item \textbf{Betaxanthin is the strongest case, and its strength is coverage.} yeast-GEM
  reaches $19\,\%$ of the screen and misses $73\,\%$ of its high producers. On genes a
  flux-based method has no features for, CGT's top 25 is enriched $4.4\times$ over base
  rate. Where both methods can be scored they are tied, they find largely different genes
  (9 shared of 29 each), and their union reaches 49 of 100.

\item \textbf{Configuration choice moves results more than method choice.} Betaxanthin grid
  cells span Spearman $0.013$ to $0.158$ against a between-method gap of $0.0014$ on AUC.
  The amino-acid strand's best run is $0.211$ and its median is $0.0498$.
\end{enumerate}

\subsection{The decision\secstatus{todo}}

\begin{itemize}
\item \textbf{Figure 6 is ready to be argued} on coverage rather than on accuracy, and the
  amino-acid material belongs in it as a coupling measurement, not as a joint-training
  result.
\item \textbf{Figure 3 is not ready, and its placeholder numbers are far from the measured
  ones.} A per-gene Pearson of $0.24$ against a ceiling of $0.775$ is not a headline. The
  cheapest path to a better one is not another architecture ladder: it is training
  morphology on its full dataset, finding where the expression curve actually turns, and
  checking whether the graph prior is the right prior at all.
\item \textbf{The graph-prior probe has run and it closed a planned arm.} The
  network-distance pair term is not worth building, the nine masked heads should be freed,
  and the mask should stop symmetrizing its two directed relations. Two cheap experiments
  remain unrun: morphology at full scale, and one expression run trained until the
  validation curve turns.
\end{itemize}

%% notes-tex/019-simb-multimodal/sections/2-expression.tex
%% [[experiments.019-simb-multimodal.expression-round-retrospective]]
%% SOURCE: numbers are from experiments/019-simb-multimodal/results/*.json (committed
%% artifacts) and from named W&B run ids in project zhao-group/torchcell_019_expr_v8.
%% Nothing here is hand-authored.

\section{Expression\secstatus{todo}}
\label{sec:expression}

\subsection{The task and its ceiling\secstatus{todo}}

The \term{expression strand} predicts, for a strain carrying one gene deletion, the
$\log_2$ ratio of every measured gene's transcript abundance against wild type. Targets
come from the Kemmeren 2014 knockout microarray compendium (doi:10.1016/j.cell.2014.02.054)
and the Sameith 2015 double-deletion compendium (doi:10.1038/sdata.2015.15), fused in the
\file{fig3_core} build. Of the 1,556 records in that build, 1,484 are single deletions and
72 are Sameith doubles, so $95.4\,\%$ of the training signal carries exactly one perturbed
gene.
\src{experiments/019-simb-multimodal/results/fig3_overlap_census.json}

The \term{reproducibility ceiling} is the highest correlation any predictor could reach
against a target that carries measurement noise. It is estimated here from the 82 deletions
measured independently in both compendia, as the mean over reporter genes of
$\sqrt{\operatorname{clip}(\hat\rho_g, 0, 1)}$ where $\hat\rho_g$ is that gene's test-retest
correlation across the paired strains. The estimate is $\mathbf{0.7746}$.
\src{experiments/019-simb-multimodal/results/expression_ceiling_replicate.json}

The best validation score observed on this task is
\file{pearson_per_feature} $= \mathbf{0.2407}$, from run \texttt{hx8pxdic} in project
\file{torchcell_019_expr_v9}, whose peak sits at epoch 9,188 of the 9,997 it reached.
The best in \file{torchcell_019_expr_v8} is $0.2293$, run \texttt{b50f93ju}, peaking at
epoch 4,007 of 4,105. Both are $\approx 30\,\%$ of the ceiling, and both peak in the last
tenth of their run.
\src{experiments/019-simb-multimodal/results/round_leaderboards.csv}

The metric is the per-reporter correlation across strains, averaged over reporters, which
is the quantity that goes to zero when a model predicts each gene's mean; the companion
per-strain metric stays near $0.4$ under exactly that collapse and is not the objective.
The scoring rule is \file{roll_max}, the maximum of a centered five-point rolling mean of
the validation curve. It is an upward-biased order statistic whose bias grows with the
number of epochs run, so it is comparable only between runs of similar length, which is why
the epoch count travels with every number quoted from the leaderboard.

The manuscript's Results section currently states $r = 0.543$ for expression. That string
sits inside a paragraph marked \texttt{[FILLER -- R3.]} and is a placeholder, not a
measurement.

\subsection{The finding of the round: nothing was trained to saturation\secstatus{todo}}

Sixty-seven runs across waves 1 to 4b were scored, and not one of them ran past 300 epochs.
Against a validation curve that dips at epochs 200 to 300 and then climbs past its early
peak, every one of those runs was stopped inside the dip.
\src{experiments/019-simb-multimodal/results/decoder_arms_torchcell_019_expr_v8.csv}

\begin{table}[htbp]
\centering\small
%% SOURCE: recomputed in the round retrospective from decoder_arms_torchcell_019_expr_v8.csv
\begin{tabular}{lrl}
\toprule
wave & runs & epoch range \\
\midrule
wave 1  & 10 & 70--140 \\
wave 2  &  5 & 76--140 \\
wave 3  & 36 & 18--276 \\
wave 4a &  8 & 77--81 \\
wave 4b &  8 & 300 \\
\midrule
\textbf{all} & \textbf{67} & \textbf{18--300, none above 300} \\
\bottomrule
\end{tabular}
\caption[]{Every scored arm in the expression round stopped at or below 300 epochs.}
\label{tab:expr-truncation}
\end{table}

Four long runs were later allowed to continue. All four died at 12.60 to 12.61 hours of
wall clock at 1,620 to 1,625 epochs, which is a scheduler limit rather than convergence,
and their validation peaks sit at 93 to $99\,\%$ of the way through the run. Eval-mode
training Pearson reached $0.652$ to $0.725$ and was still climbing when they were killed.

Longer runs have since been allowed. The picture did not change: the best v8 run peaks at
epoch 4,007 of 4,105 and the best v9 run at epoch 9,188 of 9,997. Raising the epoch budget
from 1,600 to 10,000 moved the score from $0.204$ to $0.241$ and moved the peak with it, so
the curve is still rising where it was cut off. \textbf{No peak has ever been observed on
this task}, at any budget tried, and the saturation epoch remains unknown.

\notesec{Consequence} An arm comparison at 300 epochs is not a small effect measured
noisily. Both arms were stopped before either had the chance to separate, so the difference
between them reflects early-training transients. The null sink arm reads
$+0.0024 \pm 0.0090$ over four seed pairs and post-perturbation mixing reads
$-0.0023 \pm 0.0076$, and neither number is evidence about its mechanism. More seeds do not
repair this: replicates buy precision on the wrong estimand.

\subsection{Loss and metric move in opposite directions\secstatus{todo}}

Validation \texttt{expression/loss} reaches its minimum long before validation Pearson
peaks, then rises for the remaining $\approx 1{,}200$ epochs while Pearson keeps climbing.

\begin{table}[htbp]
\centering\small
%% SOURCE: experiments/019-simb-multimodal/results/wave_history_stage-wave4b.csv
\begin{tabular}{lll}
\toprule
run & val loss min & val Pearson max \\
\midrule
\file{N_sink_mm} s1  & 0.2640 @ ep 53  & 0.0640 @ ep 257 \\
\file{N_sink_mm} s42 & 0.2653 @ ep 217 & 0.0824 @ ep 296 \\
\file{R_ref} s1       & 0.2639 @ ep 42  & 0.0632 @ ep 293 \\
\file{R_ref} s42      & 0.2662 @ ep 202 & 0.0701 @ ep 297 \\
\bottomrule
\end{tabular}
\caption[]{Minimum validation loss precedes maximum validation Pearson by 40 to 250 epochs
in wave 4b, and by $\approx 1{,}200$ epochs in the long runs. The 010 fitness and
interaction task does not do this.}
\label{tab:expr-divergence}
\end{table}

Two things follow. Every checkpoint the expression strand ever saved by validation loss was
the early-dip model, roughly 1,300 epochs before the good one, which is why
best-by-metric checkpointing was added. And the divergence itself is information about the
objective: a quantile loss on a $\log_2$ ratio is minimized by a well-calibrated, narrow
predictive distribution, while Pearson rewards getting the ordering right at whatever
scale. The prediction standard-deviation ratio moving from $0.001$ to $0.52$ over the same
window is consistent with the model widening its predictions long after the loss says it
should stop.

\notesec{What this does not license} The mechanism sentence above is a reading of the
curves, not a measurement. Nothing has tested whether a different loss recovers the
Pearson gain earlier. That is one of the options in \cref{sec:directions}.

\subsection{The perturbation operator has no pair term\secstatus{todo}}

At $|S_b| = 1$ the cross-attention that carries the perturbation runs its softmax over a
one-element key set, so the attention weight is identically 1. Re-drawing $W_Q$ and $W_K$
at standard deviation 10 changes the output by exactly $0.0$, and 16,200 of 32,760
attention parameters are dead.
\src{experiments/019-simb-multimodal/results/perturbation_selector_degeneracy.json}

The model therefore reduces to $g(h_i + c_b)$: gene identity $i$ and strain identity $b$
meet once, by addition, and no term depends on the pair $(p, i)$ of deleted gene and
reporter gene. The rank of the pair term each candidate mechanism can express is what
wave 6 was built to ladder.

\begin{table}[htbp]
\centering\small
%% SOURCE: experiments.019-simb-multimodal.multiplicative-perturbation-conditioning (derivation)
\begin{tabular}{lll}
\toprule
mechanism & form & pair rank \\
\midrule
\file{V_ref} additive & $g(h_i + c_b)$ & \textbf{0} \\
\file{V_sink} gated attention & one bounded scalar per head & 9 \\
\file{V_basis}$\{16,32,64\}$ & $\sum_j b_{ij}(h_i)\,a_j(z_S)$ & $r$ \\
\file{V_film}, \file{V_hadamard} & diagonal multiplicative & $d = 90$ \\
\bottomrule
\end{tabular}
\caption[]{Pair-rank ladder. The reference arm cannot express any dependence on which
reporter is being predicted given which gene was deleted.}
\label{tab:pair-rank}
\end{table}

Capacity is not what binds. A rank-$r$ reconstruction of the target matrix reaches $0.7265$
at $r = 32$ and $0.7799$ at $r = 64$, against a best model score of $0.241$, so a null on
this axis is a statement about the mechanism.
\src{experiments/019-simb-multimodal/results/lowrank_output_ceiling.json}

\subsection{The largest signal in the strand is imputation, and it is orthogonal\secstatus{todo}}

A ridge oracle that observes $m$ randomly chosen genes of a held-out strain and predicts
the rest reaches far past anything the model does.

\begin{table}[htbp]
\centering\small
%% SOURCE: experiments/019-simb-multimodal/results/masked_conditioning_oracle.json
%% SOURCE: experiments/019-simb-multimodal/results/cross_study_conditioning_oracle.json
\begin{tabular}{lrrrr}
\toprule
$m$ observed & within-Kemmeren & within-Sameith & Kem $\to$ Sam & Sam $\to$ Kem \\
\midrule
10   & 0.4562 & 0.3649 & 0.2335 & 0.2122 \\
100  & 0.6693 & 0.6417 & 0.3815 & 0.3922 \\
1000 & 0.7832 & 0.7779 & \textbf{0.4838} & \textbf{0.4803} \\
\bottomrule
\end{tabular}
\caption[]{Masked-conditioning oracle, \file{pearson_per_feature} on validation
strains. The within-study value at $m = 1000$ exceeds the $0.775$ reproducibility ceiling,
which the cross-study column resolves: same-array conditioning was predicting the target's
own measurement noise.}
\label{tab:oracle}
\end{table}

The gain is not a better genotype-to-expression map. Removing a genotype predictor, a
$k$-nearest-neighbor model on \file{prot_T5_all} with $k = 25$, leaves $97.5$ to
$100.6\,\%$ of the conditioning gain intact, so the channel the objective optimizes is
switched off at $m = 0$, which is where the strand is scored.
\src{experiments/019-simb-multimodal/results/conditioning_gain_after_genotype.json}

\tcfigfit{figures/residual_covariance_diagnostic.pdf}{Residual gene-gene structure in the
expression compendium after removing each gene's mean. The pattern replicates on held-out
strains and lives in about 33 effective components, which is what a per-gene independent
readout discards.}{fig:residual-cov}

What makes the oracle worth keeping is the structure it exploits. After removing each
gene's mean, the residual gene-gene correlation pattern replicates on held-out strains at
split-half $r = 0.8687$ against a permutation null of $8.45\times10^{-5}$, and it has an
effective rank of $32.78$, with $59.1\,\%$ of variance inside the top 32 components. A
per-gene independent readout emits 6,127 marginals and discards all of it.
\src{experiments/019-simb-multimodal/results/residual_covariance_diagnostic.json}

\subsection{Why the byte count is misleading\secstatus{todo}}

The expression compendium holds more scalar labels than the trigenic interaction data the
model does well on, which is the intuition that predicted this task would be easy. The
label count is not the axis that matters.

\begin{table}[htbp]
\centering\small
%% SOURCE: row counts read from $DATA_ROOT/data/torchcell/<dataset>/preprocess/data.csv
\begin{tabular}{lrrrp{42mm}}
\toprule
dataset & records & per record & scalars & perturbation axis \\
\midrule
Kuzmin 2018 trigenic interaction  & 91{,}111 & 1       & 91{,}111    & triples; a gene recurs with many partners \\
Kemmeren 2014 knockout expression & 1{,}484  & 6{,}169 & 9{,}153{,}196 & 1{,}484 single deletions, each once \\
Ohya 2005 morphology              & 4{,}718  & 281     & 1{,}325{,}758 & 4{,}718 single deletions, each once \\
\bottomrule
\end{tabular}
\caption[]{Scalar labels are not the binding resource. Expression carries $\approx 9.2$
million scalars against the trigenic set's 91 thousand, but its perturbation axis has 1,484
points and no gene is deleted twice, so nothing in the data isolates the effect of a
deletion from the effect of the one array it was measured on.}
\label{tab:label-accounting}
\end{table}

In the trigenic data a given gene appears in many records paired with different partners,
so the supervision repeatedly constrains that gene's contribution. In the expression data a
gene appears as the perturbation exactly once. The 6,169 labels attached to it are 6,169
readouts of a single perturbation event, correlated with each other at the level the
residual covariance above measures, so they are worth far less than 6,169 independent
constraints on the deletion.

\notesec{Established} The expression strand sits at $0.241$ against a ceiling of $0.775$.
The round produced no valid comparison between mechanisms, because every arm was scored
before saturation. It produced three things that hold: the additive operator provably
cannot express a pair term; loss and metric diverge on this target and not on 010; and the
one large available signal is imputation from other genes of the same strain, measured to
be orthogonal to the genotype-only task.

\notesec{Cannot say} Whether any of the pair-rank mechanisms help. Whether the graph prior
is the right prior for this task. Whether a different objective closes the loss-metric gap.

%% notes-tex/019-simb-multimodal/sections/3-morphology.tex
%% [[experiments.019-simb-multimodal.scripts.calmorph_variance_analysis]]
%% SOURCE: ceiling + shortlist from
%% experiments/019-simb-multimodal/scripts/morphology_noise_ceiling.py ->
%% results/morphology_noise_ceiling.json and results/morphology_feature_ceiling.csv.
%% Scores from results/round_leaderboards.csv.

\section{Morphology\secstatus{todo}}
\label{sec:morphology}

\subsection{The target\secstatus{todo}}

The \term{morphology strand} predicts the CalMorph phenotype vector of a single-deletion
strain. CalMorph's nominal vocabulary is 501 parameters, which resolve into 281 base
parameters and 220 coefficient-of-variation statistics; the served target is the 281 base
parameters, of which the model scores 278 after three near-degenerate features are dropped.
Source is Ohya 2005 (doi:10.1073/pnas.0509436102), 4,718 deletion mutants plus 122
independent \gene{his3} wild-type replicates.

Raw CalMorph values span about eight orders of magnitude, from $6.7\times10^{-5}$ to
$1.49\times10^{4}$ in absolute mean, so an unnormalized mean-squared-error loss is
dominated by whole-cell size and actin-region size and read $\approx 4.7$ million. The
training path now applies a per-feature $z$-score fit on the training split only, which
brings the loss to order 1 and makes ``predict each feature's mean'' score zero rather than
win. Ohya's own normalization is a per-parameter Box-Cox fit on wild-type data followed by
standardization, and it discards 247 of 501 parameters on a Shapiro-Wilk normality test at
$p \ge 0.5$. That verdict is a normality test over base and CV parameters together, not a
variance floor on the 278 modeled here, so it was not imported.

\subsection{The ceiling was never measured until now, and it is high\secstatus{todo}}

Ohya measures each of the 4,718 deletion strains once and the \gene{his3} wild type 122
independent times, which is exactly what makes reliability estimable without replicating the
mutants. Each strain's value for feature $k$ is one noisy draw $T_k = S_k + E_k$: the
strain's true value plus one sample's measurement noise. The wild-type replicates hold
genotype fixed, so their spread is $E_k$ alone; the mutant panel holds nothing fixed, so its
spread carries $S_k$ and $E_k$ together.

\notesec{Reliability and ceiling are the same measurement on two scales, which is why both
names exist} \term{Reliability} is the fraction of observed across-strain \emph{variance}
that is real signal,
$$\rho_k = \frac{\operatorname{Var}(S_k)}{\operatorname{Var}(T_k)}
  = 1 - \frac{\sigma^2_{\mathrm{WT},k}}{\sigma^2_{\mathrm{mutants},k}},$$
estimated with the sample variance over the 122 replicates and over the 4,718 mutants. The
\term{ceiling} is what a \emph{correlation} can reach. A perfect predictor outputs $S_k$ and
is still scored against the noisy $T_k$, so it correlates at
$$\operatorname{corr}(S_k,\, S_k + E_k)
  = \frac{\operatorname{Var}(S_k)}{\sqrt{\operatorname{Var}(S_k)}\sqrt{\operatorname{Var}(T_k)}}
  = \sqrt{\rho_k} = c_k.$$
The square root is the entire difference: correlation is a ratio of standard deviations
where reliability is a ratio of variances. Since $\rho_k \le 1$ the ceiling always exceeds
the reliability, so the two must never be quoted interchangeably. Over the 278 modeled
features the means are $\rho = 0.444$ and $c = 0.611$, and those are one measurement stated
twice. Compare the ceiling to a reported Pearson; compare the reliability to an $R^2$.

The \term{robust CV} used alongside them is a different axis entirely, measuring how far
deletions move a feature rather than how well it is measured:
$\mathrm{rCV}_k = (Q_{75} - Q_{25}) / |\mathrm{med}|$ over the 4,718 mutants, using
quartiles rather than a standard deviation because the CalMorph features are bounded ratios
and counts with heavy tails.

\begin{table}[htbp]
\centering\small
%% SOURCE: experiments/019-simb-multimodal/results/morphology_noise_ceiling.json
\begin{tabular}{lr}
\toprule
quantity & value \\
\midrule
mutants & 4{,}718 \\
wild-type replicates & 122 \\
modeled features & 278 \\
mean reliability & 0.444 \\
\textbf{mean ceiling over modeled features} & \textbf{0.611} \\
median ceiling & 0.647 \\
features with ceiling $> 0.5$ & 201 of 278 ($72\,\%$) \\
\midrule
best observed (\texttt{vsceij2v}, peak epoch 27 of 65) & \textbf{0.0824} \\
fraction of ceiling realized & $\mathbf{13.5\,\%}$ \\
\bottomrule
\end{tabular}
\caption[]{Morphology has the second-highest ceiling of any strand here and the lowest
fraction of it realized. Nearly three quarters of the modeled features are individually
measurable to $r > 0.5$.}
\label{tab:morph-ceiling}
\end{table}

\subsection{The strand was never really run, on two counts\secstatus{todo}}

Twenty-three runs exist in \file{torchcell_019_morph_v5}, spanning 11 to 121 epochs,
one of them a collapsed constant predictor. The best peaks at epoch 27 of 65. The joint
expression-plus-morphology project is thinner still: 32 runs, none past 87 epochs, best
\file{val/mean/pearson_per_feature} of $0.0505$.

Set against \cref{sec:expression}, where the same architecture on a related target was
still improving at epoch 9,697, a 65-epoch morphology run carries no information about
whether morphology is learnable. The morphology result is not a null. It is an absence of
measurement, and the two must not be recorded the same way.

\notesec{And the training set was a quarter of what exists, in every project} This was
checked against all eight projects that carry a morphology head rather than against one,
because the first version of this claim rested on \file{torchcell_019_morph_v5} alone and
the natural objection is that some other round did train on the full screen. None did.

\begin{table}[htbp]
\centering\small
%% SOURCE: results/round_leaderboards.csv, n_train_supervised per project
\begin{tabular}{lrrr}
\toprule
group & projects & runs & $n_{\mathrm{train}}$ observed \\
\midrule
morphology alone (\file{morph_v2}, \file{_v3}, \file{_v5}) & 3 & 183 & 1{,}161 \\
morphology with expression (\file{expr_morph} .. \file{_v5}) & 5 & 214 & 1{,}161 \\
\midrule
\textbf{total} & \textbf{8} & \textbf{397} & \textbf{1{,}161, no other value} \\
\bottomrule
\end{tabular}
\caption[]{Every morphology run ever launched, across all eight projects that carry a
morphology head, trained on the same 1,161 strains. The Ohya screen holds 4,718 deletions
and 1,440 of them also carry expression; the configs set
\file{require_modalities: [expression_log2_ratio, calmorph]}, which pins every arm to that
intersection, and the 80/20 split takes it to 1,161. For contrast, the same census shows the
expression strand's training set growing across rounds (1,125 to 1,253) and the betaxanthin
rounds differing (3,694 to 4,235), so a constant is a finding here rather than an artifact
of how the census reads.}
\label{tab:morph-ntrain}
\end{table}
\src{experiments/019-simb-multimodal/results/project_census.json}
\src{experiments/019-simb-multimodal/conf/delta_joint_expr_morph_000.yaml}
\src{experiments/019-simb-multimodal/results/fig3_overlap_census.json}

The claim does not rest on the table alone. \file{require_modalities} is set in the config
that every morphology arm inherits, and the overlap census independently gives
\file{ohya_with_expression} $= 1{,}440$ against \file{ohya_unique} $= 4{,}718$. The W\&B
column is corroboration of a restriction that is visible in the configuration.

That restriction is correct \emph{for the auxiliary-task question}, since comparing a joint
model against a morphology-only model on different instance sets would confound the
auxiliary head with a data-quantity difference. It is wrong for the question ``how well can
morphology be predicted'', and the same configs were used for both. \textbf{Morphology has
never been trained on its own full dataset.}

\subsection{Which second label should morphology be paired with\secstatus{todo}}

The overlap census answers this directly, and the answer is not expression.

\begin{table}[htbp]
\centering\small
%% SOURCE: experiments/019-simb-multimodal/results/fig3_overlap_census.json, db_overlap_census
\begin{tabular}{lr}
\toprule
genotypes carrying & count \\
\midrule
morphology (Ohya, total) & 4{,}718 \\
morphology and fitness & \textbf{4{,}220} \\
morphology and expression & 1{,}440 \\
morphology, fitness and expression & 1{,}326 \\
\bottomrule
\end{tabular}
\caption[]{Fitness shares nearly three times as many strains with morphology as expression
does, and covers $89\,\%$ of the morphology screen against expression's $31\,\%$.}
\label{tab:morph-overlap}
\end{table}

Pairing morphology with fitness keeps 4,220 of the 4,718 strains, so the joint arm and the
morphology-only arm can be run on a large common instance set without the restriction above
costing three quarters of the data. Pairing it with expression caps both arms at 1,440.

\notesec{The cap is a consequence of \file{require_modalities}, not a fact about the data}
Both caps come from dropping any genotype that is missing either phenotype. Masking the loss
over absent outputs instead lets a strain with morphology and no expression still contribute
morphology gradient, and a strain with both contribute both. That removes the cap on the
\emph{training} side while leaving the controlled comparison available on the shared subset,
and it is the documented way around this. No run here did it; \file{require_modalities} is
set in the config every morphology arm inherits.

\notesec{The question to ask, stated as one sentence} \textbf{Does adding fitness as a second
label improve morphology prediction at all, on any CalMorph feature or on the average across
them?} That is a yes-or-no question with a correct control (the same architecture, the same
strains, morphology head only), and it is worth answering before any more elaborate
multi-task claim. Report it per feature as well as averaged, because a second label that
helps the 201 well-measured features and does nothing for the rest would be invisible in the
mean.

\notesec{This does not retire the expression-helps-morphology question} That question is
what \file{delta_joint_expr_morph_000.yaml} was built to answer, its control is right, and
its three conditions (\texttt{expr}, \texttt{morph}, \file{expr_morph}) are the correct
decomposition. It has simply never been run past 87 epochs. The intended design is the
broader one: train on all of morphology with masked outputs, then ask whether adding
expression supervision on the 1,440 strains that carry it improves anything. That is a long
and awkward run to tune, which is the practical reason expression is being settled first.

\notesec{What the earlier ``0.05 to 0.14 rise'' was} A flatten-Pearson computed over the
$[B, 281]$ matrix before per-feature normalization, which is dominated by between-feature
scale: large features are large in both prediction and target. After $z$-scoring, the same
smoke run reports $\approx 0$ on that metric. The honest metric is the per-feature
correlation across strains, averaged over features, which is what
\cref{tab:morph-ceiling} scores.

\subsection{A scalar target is a reasonable next step, and this is which scalar\secstatus{todo}}

Ranking single features on reliability alone selects ones that are precisely measured and
nearly constant across mutants, which is the wrong end of the trade. The rank key is
$s_k = c_k \cdot \mathrm{rCV}_k$, the ceiling multiplied by the robust coefficient of
variation, so a feature has to be both measurable and actually moved by deletions.

%% GENERATED by experiments/019-simb-multimodal/scripts/morphology_noise_ceiling.py
%% SOURCE: experiments/019-simb-multimodal/results/morphology_noise_ceiling.json, scalar_shortlist
%% Descriptions are verbatim from Ohya 2005 SI Table 2 "parameter statistics"
%%   $DATA_ROOT/torchcell-library/ohyaHighdimensionalLargescalePhenotyping2005/si/si6.pdf
%%   sha256 037a50927bc5b978ab38b207b264da87550feeeeaaf21bc311f65a5ef13e9638
\begin{table}[htbp]
\centering\footnotesize
%% Ragged-right paragraph columns: a narrow cell holding one long underscored token
%% cannot be justified, and justifying it emits an underfull hbox on every row.
\begin{tabular}{l >{\raggedright\arraybackslash}p{44mm} >{\raggedright\arraybackslash}p{72mm} r r r}
\toprule
feature & CalMorph label & SI description & ceiling & robust CV & score \\
\midrule
\file{A113} & \file{actin_n_ratio} & No actin patch found ratio & $0.873$ & $2.369$ & $2.068$ \\
\file{C124_C} & \file{Medium_bud_ratio} & Ratio of medium bud on stage C & $0.698$ & $0.886$ & $0.618$ \\
\file{A110_A1B} & \file{Actin_f_ratio} & Actin f ratio on stage A1B & $0.637$ & $0.939$ & $0.598$ \\
\file{C125_A1B} & \file{Large_bud_ratio} & Ratio of large bud on stage A1B & $0.784$ & $0.630$ & $0.494$ \\
\file{A105} & \file{actin_a_ratio} & Actin a ratio & $0.655$ & $0.617$ & $0.404$ \\
\file{D212} & \file{nuclear_B_ratio_to_nuclear_AA1BC_cells} & Ratio of B (Nuclear) to A, A1, B and C cells & $0.604$ & $0.642$ & $0.387$ \\
\file{D201} & \file{nuclear_B_ratio} & Ratio of B (Nuclear) & $0.565$ & $0.641$ & $0.362$ \\
\file{A109_A1B} & \file{Actin_e_ratio} & Actin e ratio on stage A1B & $0.476$ & $0.756$ & $0.360$ \\
\file{A103_A1B} & \file{Relative_distance_of_actin_patch_center_from_neck_in_mother} & Relative Distance of actin patch center from neck in mother cell on stage A1B & $0.812$ & $0.431$ & $0.350$ \\
\file{A103_C} & \file{Relative_distance_of_actin_patch_center_from_neck_in_mother} & Relative Distance of actin patch center from neck in mother cell on stage C & $0.534$ & $0.639$ & $0.342$ \\
\file{A110} & \file{actin_f_ratio} & Actin f ratio & $0.700$ & $0.468$ & $0.328$ \\
\file{D215} & \file{nuclear_B_ratio_to_nuclear_A1BC_cells} & Ratio of B (Nuclear) to A1, B and C cells & $0.480$ & $0.578$ & $0.277$ \\
\file{A8-2_C} & \file{Total_brightness_of_actin_region_in_bud} & Actin region brightness in bud on stage C & $0.594$ & $0.449$ & $0.267$ \\
\file{A8-2_A1B} & \file{Total_brightness_of_actin_region_in_bud} & Actin region brightness in bud on stage A1B & $0.595$ & $0.436$ & $0.259$ \\
\file{A121_A} & \file{Maximal_distance_between_patches} & Maximam actin patch length on stage A & $0.492$ & $0.517$ & $0.254$ \\
\file{A8-1_A} & \file{Actin_region_brightness} & Actin region brightness in mother cell on stage A & $0.562$ & $0.444$ & $0.250$ \\
\file{A120_A} & \file{Total_length_of_actin_patch_link} & Total length of actin patch link on stage A & $0.404$ & $0.613$ & $0.248$ \\
\file{D208} & \file{nuclear_B_ratio_to_budded_cells} & Ratio of B (Nuclear) to budded cells & $0.424$ & $0.575$ & $0.244$ \\
\file{A108} & \file{actin_d_iso_ratio} & Actin d (iso) ratio & $0.641$ & $0.375$ & $0.241$ \\
\file{C122} & \file{large_bud_ratio} & Ratio of large bud & $0.693$ & $0.342$ & $0.237$ \\
\bottomrule
\end{tabular}
\caption[]{The twenty best scalar morphology targets, ranked by ceiling multiplied by
robust CV. The list is dominated by bounded class ratios, actin localization class,
bud-size class and nuclear-stage class, which is what a deletion actually moves. Cell
size is the opposite trade and cell shape more so: cell size at the unbudded stage
\file{C11-1_A} has a ceiling of $0.928$ but a robust CV of $0.097$, ranking 101st, and
mother-cell roundness \file{C115_A} has a ceiling of $0.972$ with a robust CV of $0.024$,
ranking 245th of the 275 features with a defined rank key. That last id was described as a
whole-cell axis ratio in the previous revision; SI Table 2 calls it roundness of the mother
cell, and it is a shape parameter rather than a size one. Descriptions are quoted
verbatim from Ohya 2005 SI Table 2, source typography included, and the label column is
the same parameter's name in SI Table 1; the two tables word some parameters
differently, and both wordings are the paper's. In a CalMorph id the
leading letter names the stain the parameter is measured from, C the FITC-ConA cell-wall
image, A the rhodamine phalloidin actin image and D the DAPI nuclear image, and the
suffix names the nuclear stage the cells were binned into, \texttt{\_A}, \texttt{\_A1B}
or \texttt{\_C}, with no suffix meaning all stages pooled. The single-letter classes
inside a description, actin $a$ to $f$ and nuclear $A$ to $F$, are defined by the
drawings in SI Fig.~5B rather than by any text, so those rows rest on that figure.}
\label{tab:morph-scalar}
\end{table}


\notesec{On ``predict cell size first'', which the data does not support} The idea is
attractive because cell size is the CalMorph quantity closest to something intuitive, and it
is the one place the panel gives a clean answer. Cell size is measured superbly and moved
almost not at all.

\begin{table}[htbp]
\centering\small
%% SOURCE: experiments/019-simb-multimodal/results/morphology_noise_ceiling.json, size_features
\begin{tabular}{llrrr}
\toprule
feature & CalMorph description & ceiling & robust CV & rank of 275 \\
\midrule
\file{C11-2_A1B} & Area of daughter cell on stage A1B     & 0.880 & 0.125 & 78 \\
\file{C11-1_A}   & Mother cell size on stage A            & 0.928 & 0.097 & 101 \\
\file{C101_A1B}  & Cell size on stage A1B                 & 0.928 & 0.088 & 118 \\
\file{C103_A}    & Long axis length of mother on stage A  & 0.949 & 0.047 & 209 \\
\file{C115_A}    & Roundness of mother cell on stage A    & 0.972 & 0.024 & 245 \\
\midrule
\file{A113}      & (the shortlist leader, for contrast)   & 0.873 & 2.369 & 1 \\
\bottomrule
\end{tabular}
\caption[]{Cell-size features are the opposite trade from the shortlist: near-perfect
ceilings of $0.88$ to $0.97$, and among the least variable features in the panel. Whole-cell
area at stage A has a robust CV of $0.097$ against a population median of $0.125$, so
\textbf{173 of the 278 features have more spread than whole-cell area}, and its rank key is
23 times smaller than \file{A113}'s. \emph{Correction to revision 2}, which called
\file{C115_A} a whole-cell axis ratio: the CalMorph SI defines it as roundness of the mother
cell, a shape parameter rather than a size one. The one size feature worth noting is bud
area (\file{C11-2_A1B}), which carries $28\,\%$ more spread than mother area, plausibly
because a deletion that stalls the cell cycle changes bud size at fixed nuclear stage.}
\label{tab:morph-size}
\end{table}

\notesec{And on ``predict colony size first''} Colony size is fitness, and fitness is already
the strongest strand in the project, trained on the Costanzo and Kuzmin single-mutant fitness
data. Using it as the morphology warm-up would test the pipeline rather than the phenotype.

\textbf{Recommendation: run the scalar warm-up on \file{A113} or \file{C124_C}}, both
reliable and genuinely variable, and treat the run as a convergence-budget calibration for
the 278-feature target rather than as a result. This is queued behind the expression work
(\cref{sec:directions}) rather than proposed for now.

\begin{keybox}
\textbf{Morphology is the largest unexploited target here, and it has never been measured
properly.} Its ceiling, estimated for the first time, is $0.611$ averaged over 278 features,
with 201 of them individually above $0.5$, which is the second-highest ceiling of any strand.
The best score achieved is $0.0824$, at $13.5\,\%$ of it. But that score comes from a
1,161-strain training set drawn from a 4,718-strain screen, in a run that stopped at epoch
65, and \textbf{every one of the 397 morphology runs across eight projects trained on the
same 1,161 strains} because \file{require_modalities} pins them to the intersection with
expression. \textbf{So $0.0824$ is not a null result about morphology; it is the absence of a
measurement}, and the two must not be recorded the same way. For a second label, fitness
shares 4,220 strains with morphology against expression's 1,440.
\end{keybox}

\notesec{Cannot say} Anything about whether morphology is learnable from genotype, because
no morphology run has passed 121 epochs or seen more than a quarter of the data. Anything
about whether expression and morphology share structure, because the joint project's
longest run is 87 epochs.

%% notes-tex/019-simb-multimodal/sections/4-betaxanthin.tex
%% [[experiments.020-cachera-betaxanthin.merzbacher-comparison]]
%% SOURCE: every number in this section is written by
%% experiments/020-cachera-betaxanthin/scripts/plot_merzbacher_comparison.py into
%% experiments/020-cachera-betaxanthin/results/merzbacher_comparison_figures.json, or by
%% experiments/019-simb-multimodal/scripts/pigment_noise_ceiling.py.

\section{Betaxanthin production\secstatus{todo}}
\label{sec:betaxanthin}

\subsection{The task, its ceiling, and the two names that get confused\secstatus{todo}}

\term{Cachera} is a genome-wide betaxanthin deletion screen
(doi:10.1093/nar/gkad656): the four-gene betalain cassette \gene{CYP76AD1}, \gene{DOD},
\gene{ARO4}$^{\mathrm{K229L}}$ and \gene{ARO7}$^{\mathrm{G141S}}$ was transferred by CRI-SPA
into each strain of the yeast knockout collection and per-colony fluorescence read as a
production proxy. It supplies 4,735 single deletions and is the ground truth here.
\term{Flux Cone Learning} (FCL) is Merzbacher et al.'s method, which learns from flux
samples drawn through a genome-scale metabolic model;
\file{RandomForestClassifier_Resampled} is its best variant. The comparison is CGT
against FCL, both scored on the Cachera screen.

Betaxanthin is the only phenotype in this document with a high ceiling. Its records carry a
per-strain standard error over a median of 15 colonies, giving a reliability of $0.836$ and
a Pearson ceiling of $\mathbf{0.914}$.
\src{experiments/019-simb-multimodal/results/pigment_noise_ceiling.json}

\tcfigfit[0.62]{figures/betaxanthin_noise_ceiling.pdf}{Reliability of the betaxanthin
target. Unlike every other phenotype here, Cachera releases a per-strain standard error over
a median of 15 colonies, so the noise term is measured directly rather than inferred from
replicate structure. That is what puts this ceiling at $0.914$ and makes betaxanthin the one
strand where a regression target is clearly
supportable.}{fig:bx-ceiling}

\subsection{Three betaxanthin numbers, two scoring rules, and one honest comparison\secstatus{todo}}
\label{sec:bx-two-numbers}

\begin{table}[htbp]
\centering\small
%% SOURCE: results/analysis_002_noise_and_lambda.json; results/round_leaderboards.csv;
%% OBJECTIVE_METRIC in experiments/020-cachera-betaxanthin/scripts/optuna_metabolism_sweep.py
\begin{tabular}{llrrl}
\toprule
number & where it comes from & scoring rule & $n_{\mathrm{train}}$ & split \\
\midrule
$\mathbf{0.4301}$ & Optuna study \file{betaxanthin_002} & raw per-epoch max & 4{,}235 & ordinary ratio split \\
$\mathbf{0.4015}$ & the same runs, in W\&B & \texttt{roll\_max} & 4{,}235 & ordinary ratio split \\
$\mathbf{0.3705}$ & \file{torchcell_020_betaxanthin_v4} & \texttt{roll\_max} & 3{,}698 & Merzbacher test genes pinned out \\
\bottomrule
\end{tabular}
\caption[]{The first two rows are the \textbf{same runs under two scoring rules}, which is
what made the previous revision confusing. The Optuna objective is
\file{val/betaxanthin/pearson_per_feature_max}, a running maximum over epochs; the
leaderboard's \texttt{roll\_max} first averages five neighboring epochs, and smoothing can
only lower a maximum. Only rows two and three are comparable, and their difference,
$0.4015$ against $0.3705$, is what holding out a competitor's test set costs. Against the
$0.914$ ceiling the best is $47\,\%$ realized on the raw rule and $44\,\%$ on the smoothed
one.}
\label{tab:bx-two-numbers}
\end{table}

\notesec{On whether $0.4301$ is a plateau or one lucky trial, since the previous phrasing of
this was unreadable} A \term{trial} here is one Optuna trial, meaning one hyperparameter
configuration trained once. It is not a data split: the split is pinned across the whole
study, so trials differ by hyperparameters and by weight initialization only. Study
\file{betaxanthin_002} ran 37 of them. Two things follow, and they point in opposite
directions.

\begin{itemize}
\item \textbf{The top is flat, so the winner is not a spike.} The best five trials read
  $0.4301$, $0.4234$, $0.4216$, $0.4211$ and $0.4095$, a spread of $0.0206$. Re-running an
  \emph{identical} configuration moves the objective by $\sigma = 0.0302$
  (\cref{tab:replicate-noise}), so those five sit inside two thirds of one standard
  deviation of each other. They are statistically indistinguishable, which means the winner
  is one of five equivalent configurations rather than a single fortunate run.
\item \textbf{But taking the maximum of 37 noisy trials is still biased upward.} That is
  \term{selection inflation}, estimated at $0.064$ here. So a configuration chosen in
  advance, without seeing the study, should be expected to score around $0.366$ rather than
  $0.430$.
\end{itemize}

Both are true at once: the plateau says the number is not a fluke of one run, and the
inflation says it is not what a fresh configuration would give.

\notesec{The comparison number is lower than the best number on purpose} The head-to-head
below runs on \file{020_v4}, whose split \textbf{pins the 640 Merzbacher test genes out of
training}. That costs 537 training strains, and it costs score. Using the higher-scoring
model instead would mean scoring on genes it had trained on, which is the one thing that
would void the comparison. The cell carried forward is \file{s09_L6_maskon_lr0.0001_yj},
selected by validation at $0.3639$ and never by its score on the comparison genes.

\notesec{The 537 strains are not recoverable, and the fix is a reporting change that needs
no retraining} The natural proposal is to train on everything outside the comparison list,
carve validation and test from that same remainder, and evaluate twice. Sized against the
data, that proposal is numerically the split \file{020_v4} already ran: it gives
$3{,}682/299/295$ plus the 639 comparison genes, against \file{020_v4}'s actual
$3{,}698/286/931$, where the 931 is 639 comparison genes plus 292 of our own. Holding the
comparison genes out of training is what costs the 537, under any scheme, and no
rearrangement returns them.

\begin{table}[htbp]
\centering\small
%% SOURCE: experiments/020-cachera-betaxanthin/scripts/fcl_resplit_sizing_check.py
\begin{tabular}{lrrr}
\toprule
arm & train & validation & test \\
\midrule
unpinned (\file{torchcell_020_betaxanthin}, \file{_v3}) & 4{,}235 & 340 & 340 \\
pinned (\file{torchcell_020_betaxanthin_v4}) & 3{,}698 & 286 & 931 \\
the proposed re-split, sized & 3{,}682 & 299 & 295 + 639 \\
\bottomrule
\end{tabular}
\caption[]{Sizing the proposed re-split against what was run. The bookkeeping on the pinned
arm closes exactly: train $-537$, validation $-54$, test $+591 = 537 + 54$, the other 48
pinned genes having already been in the unpinned test set. The realized split ratio is
$86/7/7$ rather than the nominal $80/10/10$, because assignment is per-key stratified; the
config comment estimating ``roughly 511'' lost genes used the nominal ratio and understates
the true 537 by 26.}
\label{tab:bx-resplit}
\end{table}

\textbf{What the proposal does buy, and it is worth taking, is the second evaluation.} The
existing 931-record test set is already 639 comparison genes plus 292 of our own. Splitting
it and reporting both numbers separates ``how we do against FCL on their genes'' from ``how
the model does on a clean held-out sample'', and it runs off the existing \file{020_v4}
checkpoints with no retraining at all. Enlarging our own test set beyond those 292 is the
only part that costs anything, at roughly a further 137 training strains for a 429-gene own
test set.

\notesec{One name-resolution defect, now triangulated and resolved} Of the 640 comparison
genes, 630 carry a betaxanthin measurement, and the pin file holds 639. The screen reports
one gene as \gene{PPA1}, an ambiguous alias naming both \gene{IPP1} (\gene{YBR011C}), which
is essential, and \gene{VMA16} (\gene{YHR026W}), which is not. Neither owns \gene{PPA1} as
its standard name, so the resolver's standard-name-before-alias rule cannot break the tie and
correctly returns \texttt{AMBIGUOUS}. The raw screen does not settle it either: its 27
columns carry a gene name and colony statistics, with no systematic ORF, plate, position or
barcode.

Five independent lines were run, and they agree.
\src{experiments/020-cachera-betaxanthin/scripts/triangulate\_ppa1\_alias.py}

\begin{itemize}
\item \textbf{The strain is alive, which is decisive on its own.} The \gene{PPA1} row records
  \textbf{15 colonies at 24 h and 15 at 48 h}, with mean areas of $106.5$ and $118.5$. A
  haploid \gene{ipp1}$\Delta$ is inviable and yields no colonies, so no row for it can exist.
  This is the strain's own growth data and does not depend on any claim about which
  collection was screened. The colony is small, in the $1.6$th percentile of screen area,
  which is what a severe but survivable defect looks like.
\item \textbf{The screen contains essentially no essential genes}, on $1{,}140$
  counterexamples rather than the six previously cited. It covers $4{,}704$ of $5{,}467$
  non-essential genes ($86\,\%$) and $29$ of $1{,}140$ SGD inviable-null genes ($2.5\,\%$),
  and those 29 are autophagy genes and paralogs whose inviability is conditional, all of
  which grew.
\item \textbf{The V-ATPase family has a \gene{VMA16}-shaped hole.} Thirteen of the fifteen
  \gene{VMA} genes are in the screen; the only absentees are \gene{VMA9} and \gene{VMA16}.
  The file names genes by alias constantly (\gene{VMA1} as \gene{TFP1}, \gene{VMA3} as
  \gene{CUP5}), so an alias row for \gene{VMA16} is exactly what one expects.
\item \textbf{Its value sits on top of the V-ATPase cluster.} Betaxanthin $-1.566$ against
  \gene{VMA13} $-2.19$, \gene{VMA3} $-1.55$, \gene{VMA22} $-1.54$, \gene{VMA6} $-1.44$, on a
  screen median of $-0.006$. Soft corroboration, carrying nothing alone.
\item \textbf{Neighbor evidence is neutral here}, unlike the \gene{FEN1} case: neither
  candidate appears elsewhere in the screen under another name, so the ``one row per gene''
  argument that decides \gene{FEN1} does not apply.
\end{itemize}

\textbf{The \gene{PPA1} row is \gene{YHR026W} (\gene{VMA16}), and the split builder is
right.} Nothing points to \gene{IPP1} except the build itself.

\notesec{The training build has it wrong, and the head-to-head does not} The defect is real
but narrower than feared. In both the loader LMDB and the \file{fig6_pigment_transfer}
training build, \gene{YBR011C} carries the betaxanthin value $-1.5664$ and \gene{YHR026W}
carries none. The cause is verified in code rather than inferred: the build predates the
layered resolver, and the loader then took \file{candidates[0]} from the alias map, which
returns \gene{YBR011C} for \gene{PPA1}. The same line got \gene{FEN1} right by luck.

\textbf{The comparison is not corrupted}, because \file{build_merzbacher_split.py} applies
the screen-local disambiguation when it reads the raw file, so the comparison's ground truth
for \gene{YHR026W} is the \gene{PPA1} value and \gene{YBR011C} is the gene dropped from 640
to 639. What remains is one false training label out of 4,735: the model was trained on a
betaxanthin value attributed to a gene whose deletion strain cannot exist, and scored on the
same measurement under its correct ORF.

\textbf{A plain rebuild does not fix this, it loses the value.} The current loader rejects
\texttt{AMBIGUOUS}, so it would drop the row and leave neither ORF labeled. The
disambiguation table lives in the split builder and needs to reach the loader, which is
issue \#195 territory and lands with the next full graph build.

\notesec{And the same alias broke Merzbacher's pipeline too} Their released 640-gene test
list contains \emph{both} \gene{YBR011C} and \gene{YHR026W}, each labeled low, so their code
expanded the alias to both candidates rather than disambiguating. It is an independent
instance of the same failure, and it is worth knowing that 5 of their 640 test genes are
SGD inviable-null in what is otherwise a non-essential list.

\subsection{The published accuracy comparison is empty for both methods\secstatus{todo}}

FCL reports gene-level accuracy $0.700$ against a majority-class rate of $0.673$, achieved
by calling $607$ of $640$ genes medium, which is $94.8\,\%$. Our absolute binning sits at
$0.681$ doing the same thing. Forced to the true class marginal, FCL drops to $0.556$ and
CGT to $0.549$.

\notesec{The structural point} Every model that finds more high producers has worse overall
accuracy, because the only way to call more highs is to stop calling everything medium.
Accuracy selects against the capability the task is about, so it should not be the primary
metric in any writeup, including ours.

\subsection{Both methods have real signal, and it lives in the tail\secstatus{todo}}

Of the top $k$ genes each method nominates, the fraction that are truly high producers,
against a base rate of $0.156$:

\begin{table}[htbp]
\centering\small
%% SOURCE: results/merzbacher_comparison_figures.json, fig3_precision_at_k
\begin{tabular}{lccc}
\toprule
model & $k = 10$ & $k = 25$ & $k = 50$ \\
\midrule
FCL RF & 8 ($5.1\times$, $p = 10^{-5}$) & 19 ($4.9\times$, $p = 8\times10^{-12}$) & 20 ($2.6\times$, $p = 10^{-5}$) \\
CGT    & \textbf{9} ($5.8\times$, $p = 4\times10^{-7}$) & 18 ($4.6\times$, $p = 10^{-10}$) & \textbf{24} ($3.1\times$, $p = 2\times10^{-8}$) \\
\bottomrule
\end{tabular}
\caption[]{Precision at $k$ for high producers. Nine of CGT's top ten are genuine high
producers. Both curves approach the base rate by $k = 150$.}
\label{tab:precision-at-k}
\end{table}

Median percentile rank by true class is flat for both methods, at 54/49/56 for FCL and
48/51/52 for CGT across low, medium and high. Read with the table above, that says the
signal is real and confined to the tail, not that there is none: neither method separates a
typical high producer from a typical low one, and both are far better than chance at the very
top of their own ranking.

\tcfigfit{figures/merzbacher_fig3_precision_at_k_2026-08-02-23-43-08.pdf}{Precision at $k$
for high producers. (a) the comparison: FCL against the CGT cell chosen by validation, with
the nine unselected grid cells drawn as a shaded envelope rather than as nine competing
lines. The base rate of $0.156$ is drawn and labeled, and it is what a random pick of $k$
genes would return. Nine of CGT's top ten are genuine high producers, and both methods reach
the base rate by $k \approx 150$, which is where the tail signal is exhausted. (b) the same
ten cells drawn individually, one color each, which is what answers whether some cells sit
below the base rate early: \textbf{three of the ten do at $k = 50$}. Read crossings at very
small $k$ with care, since at $k = 10$ precision moves in steps of $0.1$ and a crossing is a
single gene. The spread across cells is far larger than the gap between the two methods,
which is the hyperparameter sensitivity of \cref{sec:terms} rather than a split
effect.}{fig:precision-at-k}

\tcfigfit{figures/merzbacher_fig1_scatter_spread_2026-08-02-23-43-08.pdf}{Per-gene spread
over the 639 shared genes. Both models emit a nearly constant value: $80\,\%$ of FCL's
genes fall inside an expected-class band of $0.22$ on a $0$ to $2$ scale, $80\,\%$ of CGT's
inside $0.33$ $z$-units. Panel c shows the disagreement concretely, with each method
recovering 29 of the 100 true high producers and sharing 9.}{fig:bx-spread}

\subsection{The two methods find different genes\secstatus{todo}}

Rank-matched into the same 100 high slots, FCL recovers 29 of the 100 true high producers
and CGT recovers 29, and they share 9. The union reaches 49, so running both finds $69\,\%$
more high producers than running either.

\notesec{The two methods' predictions are anti-correlated, and here is what that sentence
means} Compare their predicted rankings over the 639 shared genes directly, ignoring the
truth. Two methods that are each independently right about the same target should agree
somewhat, and the amount is calculable: if one correlates with the truth at $+0.105$ and the
other at $+0.039$, then agreeing only through the truth would put their mutual correlation
at about $0.105 \times 0.039 = +0.004$, essentially zero. The observed correlation between
them is $\mathbf{-0.128}$, which at $n = 639$ is $z = -3.2$: they disagree substantially
more than two independent methods would.

\notesec{No mechanism is claimed for the negative sign} It could be a genuine
anti-relationship between what the two feature spaces encode, or an artifact of how each was
fit; nothing here distinguishes those. What \emph{is} claimed is the operational
consequence, which does not depend on the sign's explanation: because the two find largely
different genes, their recall adds rather than overlaps, so running both is worth more than
picking the better one. That fact is invisible in every aggregate statistic, all of which
report the two as tied. The 51 high producers that neither finds is the honest ceiling on
this comparison.

\subsection{The real difference is coverage\secstatus{todo}}

yeast-GEM contains 1,161 genes, of which 915 appear in the Cachera screen. That is
$19.4\,\%$ of its 4,721 deletions. A flux-based method has features for those and for
nothing else.

\begin{table}[htbp]
\centering\small
%% SOURCE: results/merzbacher_comparison_figures.json, fig8_screen_coverage
\begin{tabular}{lrr}
\toprule
population & in yeast-GEM & outside \\
\midrule
every deletion in the screen & 915 (19\,\%) & \textbf{3{,}806 (81\,\%)} \\
top 100 measured producers & 32 (32\,\%) & \textbf{68 (68\,\%)} \\
all high producers & 60 (27\,\%) & \textbf{163 (73\,\%)} \\
\bottomrule
\end{tabular}
\caption[]{Metabolic genes are enriched among high producers, at $26.9\,\%$ of the highs
against $19.4\,\%$ of the screen, a $1.4\times$ enrichment that is exactly what the biology
predicts. It is not enough to matter: restricting to metabolism buys that concentration and
costs three quarters of the hits.}
\label{tab:gem-coverage}
\end{table}

\tcfigfit{figures/merzbacher_fig8_screen_coverage_2026-08-02-23-43-08.pdf}{Measured data
only, no model and no prediction: how little of the betaxanthin screen a flux-based method
can reach.}{fig:gem-coverage}

\notesec{How the class labels are derived, since every metric below depends on it} The
Cachera screen releases a continuous production value, and FCL's task is a three-class
one, so a rule is needed to turn measurements into low, medium and high. FCL's rule:
min-max scale production onto $[0, 1]$ over their gene set and cut at $0.40$ and $0.65$.
\textbf{We apply their thresholds rather than inventing our own}, which is what makes the
two panels comparable. Two consequences follow and both are stated rather than hidden.

\begin{itemize}
\item \textbf{Applying their rule to our larger copy of the same screen does not reproduce
  their class counts.} It gives $107/476/56$ against their $109/431/100$, and $18.8\,\%$ of
  genes land in a different bin than our measured values would give. A model predicting our
  measurement \emph{perfectly} is therefore marked wrong on about a fifth of genes, so their
  labels are a headroom bound on any accuracy-style metric, ours included.
\item \textbf{They never labeled the non-GEM genes at all}, because those genes are outside
  their method's reach. We derive labels there with the identical rule, which is the only
  way that panel exists. The comparison to make is therefore panel against panel rather than
  number against number, and \cref{fig:capability-gap} is drawn that way.
\end{itemize}

\notesec{Why every metric here is rank-based or re-binned} We produce a regression and they
produce a three-class classifier, so nothing can be compared directly. Two repairs are used
throughout, and the choice between them is not arbitrary. \textbf{Rank-based metrics}
(Spearman, precision at $k$, AUC) need no thresholds at all and so cannot be affected by the
binning disagreement above; they are preferred wherever available. Where a class metric is
unavoidable, our continuous predictions are binned \textbf{with thresholds fit on the
training split only}, never on the test genes, so the binning cannot borrow information from
the evaluation. Neither repair fixes the $18.8\,\%$ label disagreement, which sits in the
ground truth rather than in either model.

On held-out test genes split by yeast-GEM membership, with labels derived the same way on
both sides, CGT's top 25 recovers 17 of 48 true highs inside the model ($9.3\times$ over
base rate, $p = 2\times10^{-15}$) and 9 of 21 outside it ($4.4\times$,
$p = 2\times10^{-5}$).

\tcfigfit{figures/merzbacher_fig9_metabolic_vs_nonmetabolic_2026-08-02-23-43-08.pdf}{CGT
predicts high betaxanthin producers among deletions a flux-based method cannot represent.
Held-out test genes split by yeast-GEM membership, identical derived labels on both sides,
circled points are CGT's top 25. Flux Cone Learning has no prediction for the right
panel.}{fig:capability-gap}
FCL's prediction for the outside set is not poor, it is undefined: flux sampling yields no
feature for a deletion the metabolic model does not contain. The FCL test set is $639$ of
$640$ inside yeast-GEM; of our other test genes only $21$ of $294$ are.

\notesec{One qualification that must travel with this claim} ``Inside'' and ``outside'' here
mean inside or outside \textbf{yeast-GEM}, the genome-scale metabolic model, not inside or
outside anyone's training set. CGT's rank correlation with measured production is
\emph{higher} on the genes yeast-GEM does not contain, $+0.270$ against $+0.124$. That is the
opposite of what one would expect if metabolic genes were the easy ones, so it is worth
testing rather than asserting: the two correlations are measured on different,
non-overlapping gene sets, and comparing two correlations requires first mapping each
through Fisher's transform, which turns a correlation into a quantity whose sampling
distribution is approximately normal with a known variance. Doing that gives $z = 2.07$,
$p = 0.039$: a real difference, but marginal, and on populations that differ in more than
their yeast-GEM membership. \textbf{The safe reading is that CGT is at least as good outside
the metabolic model as inside it}, which is all the coverage argument needs.

\subsection{Where the comparison is soft\secstatus{todo}}

\begin{itemize}
\item \textbf{The selection is asymmetric, and it favors FCL. This is now verified rather
  than suspected, and it is stronger than the previous wording.} Revision 2 said their
  validation folds ``appear to be drawn from the same 640 genes as their test list''. From
  their Zenodo release: \file{yeast_production_validation_split.csv} assigns a fold to each
  of \textbf{exactly the same 640 systematic ORFs} that
  \file{yeast_production_test_split.csv} names. The intersection is 640 and both set
  differences are empty, in five folds of 129/128/128/128/127. Their training loop then
  computes \file{train_names = set(nontest_names) - set(val_names)}, which is a no-op
  because those genes were never in the non-test pool, so every fold trained on the full
  complement and validated on a 128-gene slice of the test set. Model selection across their
  16 variants therefore ran on test data.

  Their own CHO essentiality release does \emph{not} do this: there the validation genes are
  exactly the 1,833 flagged as non-test, sharing nothing with the 457 test genes. So the
  yeast leak is a deviation from their practice inside the same release rather than a naming
  convention, which is what makes it worth stating as a fact. And the paper alone cannot
  settle it: its Methods never mention a validation set and the test count is internally
  inconsistent (659 in the main text, 649 in Table S6, 640 in the release), so only the
  released files decide it.

  Ours, by contrast, is the validation-selected cell with the comparison genes pinned out of
  training. \textbf{So where the two land level, we are level with an opponent selected on
  the test set.}
\item \textbf{Validation Pearson does not track test AUC, and this is the hyperparameter
  sensitivity again rather than a split effect.} Four of ten CGT cells beat FCL's AUC of
  $0.570$ and the best reaches $0.695$, while the validation-selected cell reads $0.557$.
  The runner-up on validation, $0.3528$ against $0.3639$, has an AUC $0.14$ higher. Every
  one of those cells sees the same split; what differs is the hyperparameters
  (\cref{sec:terms}). So the honest statement is that \textbf{our own model-selection
  criterion does not rank cells the way the comparison metric does}, which makes neither
  ``CGT's AUC is 0.695'' nor ``CGT's AUC is 0.557'' quotable on its own. Selecting on
  validation and reporting whatever AUC follows is the defensible protocol, and it is what
  was done; the spread is reported beside it rather than being selected from.
\item \textbf{Their labels are a headroom bound}, as detailed above: $18.8\,\%$ of genes bin
  differently under their rule than our measured values would give, so a model predicting
  our measurement perfectly is marked wrong on about a fifth of genes.
\item \textbf{Regression against classification.} They tried regression and abandoned it,
  in their words ``challenging with the limited number of knockouts at the high and low
  ends.'' Our ceiling of $0.914$ makes regression available to us. The comparison is
  therefore between a regressor and a three-class classifier, which is why every metric
  above is either rank-based or computed after binning our predictions with thresholds fit
  on the training split only.
\item \textbf{The build is stale} with respect to the gene-name resolver (issue \#195),
  which the split builder avoids by working from the raw screen, and which the training data
  still inherits until a rebuild.
\end{itemize}

\begin{keybox}
\textbf{The betaxanthin case is about coverage, and it does not need to beat anyone on
accuracy.} Where both methods can be scored, CGT and Flux Cone Learning are tied, both with
tail-confined signal at roughly $5\times$ enrichment in the top 25, and they find largely
different genes, so their recall adds: 29 each of 100 true high producers, sharing 9,
reaching 49 together. The asymmetry favors them, since their model was selected on its own
test set and ours on validation. \textbf{The argument that carries is the other one}: a
flux-based method has features for $19\,\%$ of the screen and misses $73\,\%$ of its high
producers, while CGT predicts high producers in the $81\,\%$ it cannot represent at
$4.4\times$ over base rate, and scores at least as well there as inside the metabolic model.
\end{keybox}

\notesec{Cannot say} That CGT beats FCL on any single aggregate. That the rank
anti-correlation means anything mechanistically. That the absolute top-100 membership is a
deployable call: the rank-matched comparison works by handing \emph{both} methods the true
number of high producers and asking which genes each puts in those 100 slots. That is the
right way to compare two methods, because neither can win by predicting more or fewer highs
than exist. But it is not available in deployment, where nobody knows how many high
producers a screen contains, so the top-100 precision measured here is an upper bound on
what the same model would deliver prospectively.

%% notes-tex/019-simb-multimodal/sections/5-amino-acid.tex
%% [[experiments.023-metabolome-betaxanthin-joint.scripts.betaxanthin_amino_acid_predictivity]]
%% SOURCE: coupling numbers from
%% experiments/023-metabolome-betaxanthin-joint/scripts/betaxanthin_amino_acid_predictivity.py
%% -> results/betaxanthin_amino_acid_predictivity.json. Scores from
%% experiments/019-simb-multimodal/results/round_leaderboards.csv. Convergence behavior from
%% experiments/022-mulleder-metabolome/results/training_length_analysis.json.

\section{Amino-acid pools, and whether they help production\secstatus{todo}}
\label{sec:amino-acid}

\subsection{The target\secstatus{todo}}

Mulleder 2016 (doi:10.1016/j.cell.2016.09.007) measured the intracellular concentration of
19 amino acids in 4,678 single-deletion strains. There is one replicate per strain and no
released standard error, so the dataset admits \textbf{no within-dataset reproducibility
ceiling}. Every correlation involving it is attenuated by an unmeasured amount, and that
qualification travels with every number below.

The best validation \file{pearson_per_feature} is $\mathbf{0.209}$, run
\texttt{24vlgope} in \file{torchcell_022_mulleder19_v4}, peaking at epoch 402 of 999.
The median across all 97 runs in the strand is $0.0570$, so the leaderboard is a long tail
under one good configuration rather than a plateau. Earlier rounds read $\le 0.025$, then
$0.171$, so the trajectory is upward across rounds.

\notesec{That median is corrected, and the correction matters for every median in this
document} Revision 2 reported $0.0498$ ``across the 31 runs''. Those 31 were not a sample of
the strand: they were the subset whose validation curves had been fetched, and the fetch
rule selects the top runs per project by final score and by epoch count
(\cref{sec:coverage}). A maximum-finder is the right instrument for a strand maximum and the
wrong one for a median, since the runs it skips are systematically the weak ones. Reading
all 97 moves the median up to $0.0570$. Maxima are unaffected by construction.

\subsection{These runs turn within their budget, and expression does not\secstatus{todo}}

Twenty-four runs in the preceding round were scored for their convergence shape. The
typical run peaks at epoch 48 to 134, continues for another 41 epochs with validation slope
turning negative while training slope stays positive, and is classified converged then
overfitting. None of them hit the epoch cap.
\src{experiments/022-mulleder-metabolome/results/training_length_analysis.json}

\notesec{Said carefully, because ``converged'' has already misled us once on this project}
What is observed is that the validation metric reaches a maximum inside the budget and then
declines, which is what makes these runs scorable. That is weaker than convergence, and the
distinction is not pedantic: the expression runs were also called converged, on curves that
dipped early and flattened, and only at several thousand epochs did it become clear that
Pearson was still climbing and squared error would come back down to a late minimum. The
same could in principle be true here beyond epoch 999, and no amino-acid run has been
extended to find out.

What can be said is the comparison, which does not depend on the word. Amino-acid runs reach
a maximum well inside a 999-epoch budget; expression runs have not reached one at 10,000.
The two use the same architecture and the same encoder, so whatever makes expression climb
for thousands of epochs is a property of the expression target or its loss rather than of the
model.

\subsection{The precursor hypothesis is wrong, and the profile hypothesis is right\secstatus{todo}}

Betalains derive from L-tyrosine (tyrosine to L-DOPA to betalamic acid), and a betaxanthin
is betalamic acid condensed with an amino acid. The mechanistic prediction is that a
deletion's free amino-acid pool carries information about how much betaxanthin that same
deletion makes. Both screens are keyed by systematic ORF and share 4,432 deletions, so the
prediction is directly testable with no model in the loop.

\notesec{What the prediction actually is, since it is easy to read as a model result and
is not one} No neural network appears anywhere in this subsection. The procedure is:

\begin{enumerate}
\item Join the two screens on systematic ORF. One row per deletion: 19 measured amino-acid
  concentrations from Mulleder, and one betaxanthin value from Cachera.
\item Take $\log$ of each concentration (they are skewed positive quantities) and
  $z$-score each of the 19 columns.
\item Fit \textbf{ridge regression} from those 19 numbers to the betaxanthin value, with
  the penalty chosen by inner cross-validation.
\item Score \textbf{out of fold}: five-fold cross-validation, so every prediction is made
  by a model that never saw that deletion, then correlate the held-out predictions against
  the measurements. Repeat over three fold shuffles and report the mean.
\end{enumerate}

So $r = 0.298$ means: \emph{knowing a deletion's 19 amino-acid concentrations lets you
predict its betaxanthin at $r = 0.298$ on deletions you did not fit on}. The comparison row
``tyrosine alone'' is the same procedure with one column instead of nineteen. It is a
statement about the two measurements, and it puts a floor under what any model sharing a
representation between these phenotypes could hope to exploit.

\begin{table}[htbp]
\centering\small
%% SOURCE: results/betaxanthin_amino_acid_predictivity.json (cross_validated_uncertainty_key)
\begin{tabular}{lrrrr}
\toprule
predictor of betaxanthin & $r$ & SE at this $n$ & SD across shuffles & $n$ \\
\midrule
tyrosine alone (the precursor)         & 0.0641 & 0.0150 & 0.0039 & 4{,}432 \\
\textbf{19 amino acids jointly}        & \textbf{0.2980} & 0.0137 & 0.0029 & 4{,}432 \\
\midrule
\multicolumn{5}{l}{\emph{on the 3{,}708 deletions with a Costanzo fitness record}} \\
single-mutant fitness alone            & 0.1508 & 0.0161 & 0.0028 & 3{,}708 \\
19 amino acids                         & 0.1763 & 0.0159 & 0.0025 & 3{,}708 \\
19 amino acids and fitness             & 0.2001 & 0.0158 & 0.0010 & 3{,}708 \\
19 amino acids, fitness regressed out of both & 0.1326 & 0.0161 & 0.0026 & 3{,}708 \\
\midrule
\multicolumn{5}{l}{\emph{on the 724 with no Costanzo record}} \\
19 amino acids                         & \textbf{0.4547} & 0.0295 & 0.0219 & 724 \\
\bottomrule
\end{tabular}
\caption[]{Cross-validated ridge, five folds, three shuffles; amino acids enter as $\log$
concentration, $z$-scored. \textbf{Two different spreads are reported because they answer
different questions and quoting only one is misleading.} The \emph{SE at this $n$} is the
Fisher standard error, $1/\sqrt{n-3}$ mapped back at the observed $r$, and it is the one to
use for ``how precisely is this correlation measured'': it is what separates the $0.298$
from the $0.064$, at $17$ standard errors. The \emph{SD across shuffles} is fold-assignment
noise on the same deletions, which is why it is an order of magnitude smaller; it says the
procedure is stable, not that the estimate is precise. A third spread, the SD across the
five folds, runs $0.025$ to $0.064$ and is an upper reference rather than a standard error,
since each fold estimates on $n/5$ deletions.}
\label{tab:aa-predictivity}
\end{table}

\notesec{Tyrosine is not the axis} Its marginal correlation with betaxanthin is
$r = -0.076$ ($p = 3.9\times10^{-7}$), which is negative, and it ranks sixth of nineteen by
absolute correlation. The strongest is methionine at $+0.140$, followed by aspartate
$+0.127$ and arginine $-0.118$. No amino acid exceeds $|r| = 0.15$. Correcting for the
betaxanthin side's known reliability of $0.836$ raises these by a factor of $1.094$, which
does not change the reading; the Mulleder side's correction is unmeasurable and would raise
them further by an unknown amount.

\notesec{The profile is the axis} The 19-dimensional pool predicts betaxanthin at
out-of-fold $r = 0.298$, roughly $4.7\times$ the precursor alone.

Revision 2 set that beside the CGT's own genotype-to-betaxanthin score and said a linear map
``recovers most of what the model recovers from genotype''. \textbf{That comparison was not
valid} and it is withdrawn: the two numbers use different splits and different scoring
rules, and one takes measured metabolites as inputs while the other takes only the genotype.
The question the comparison was reaching for is answered properly below, by a measurement
rather than by placing two numbers next to each other.

\subsection{Is the coupling redundant with what betaxanthin supervision already teaches\secstatus{todo}}
\label{sec:aa-redundancy}

This is the objection that decides whether the coupling is worth building on, and revision 2
did not address it. Both the amino-acid pool and betaxanthin are readouts of the same
deletion. So an $r$ of $0.298$ is equally consistent with two very different worlds:

\begin{enumerate}
\item[(i)] the pool carries betaxanthin-relevant information that a genotype-to-betaxanthin
  model does \emph{not} already extract, in which case a shared encoder has something to
  gain; or
\item[(ii)] both are shadows of one genotype axis the Cachera labels already teach, in which
  case no amount of architecture can beat betaxanthin-only training, and the $0.298$ is
  real but useless.
\end{enumerate}

\notesec{The measurement that separates them} Take a trained betaxanthin-only model, read
its \emph{held-out} predictions, and regress those out of both sides before asking what the
amino-acid pool still predicts. If world (ii) holds, nothing survives. Run over the ten
grid cells that dumped test predictions, on the 816 held-out deletions carrying both
measurements:

\begin{table}[htbp]
\centering\small
%% SOURCE: experiments/023-metabolome-betaxanthin-joint/scripts/bx_aa_incremental_over_model.py
%% -> results/bx_aa_incremental_over_model.json
\begin{tabular}{lrrrr}
\toprule
control model (grid cell) & $r$ of the model & $r$ of 19 AA & 19 AA given the model & corr of the two \\
\midrule
\file{s08_L6_maskon_lr0.0001} & $0.3564$ & $0.2720$ & $\mathbf{0.2179}$ & $0.2185$ \\
\file{s09_L6_maskon_lr0.0001} & $0.3305$ & $0.2720$ & $0.2262$ & $0.2050$ \\
\file{s01_L2_maskon_lr0.0001} & $0.3074$ & $0.2720$ & $0.2139$ & $0.2263$ \\
\file{s05_L2_maskoff_lr0.0001} & $0.2913$ & $0.2720$ & $0.2259$ & $0.1768$ \\
\file{s04_L2_maskoff_lr0.0001} & $0.2761$ & $0.2720$ & $0.2380$ & $0.1660$ \\
\bottomrule
\end{tabular}
\caption[]{Against the strongest available control, the amino-acid pool keeps $0.2179$ of
its $0.2720$, which is $80\,\%$, after the betaxanthin model's own held-out prediction is
removed from both sides. At $n = 816$ that is $6.5$ Fisher standard errors from zero. The
last column is the direct check on the same point: the model's predictions and the
pool's predictions correlate at only $0.22$, where near-total redundancy would need
something above $0.75$. Five further collapsed cells with $r_{\text{model}} \le 0.10$ are
omitted, since residualizing on a failed model removes nothing by construction.}
\label{tab:aa-incremental}
\end{table}

\notesec{So world (i) holds} There is betaxanthin-relevant information in the free
amino-acid pool that the genotype-to-betaxanthin map does not capture, and the joint arm's
failure to use it is therefore a statement about the auxiliary head rather than about the
signal.

\notesec{Three things this does not establish, kept separate because they are separate}
It controls for a \emph{linear} function of the model's prediction, so non-linear redundancy
is untouched. The strongest control available is $r = 0.3564$, below the strand's best
score, so redundancy against a better model than any that dumped predictions is not ruled
out. And ``the coupling is real and incremental'' is a different claim from ``an auxiliary
head is the mechanism that converts it into a gain''; the second is what the next subsection
finds against.

\notesec{Why fitness enters at all, since it is a third dataset and looks like a detour}
Both screens are deletion phenotypes and both respond to how well the strain grows. A slow
deletion can have a distorted amino-acid pool \emph{and} a distorted pigment readout with
no metabolic relationship between the two, which would produce exactly the correlation
above for an uninteresting reason. Costanzo single-mutant fitness (KanMX deletions, 30 C)
is the standard measurement of that shared nuisance, so it is regressed out of both sides
and the fit recomputed on the residuals. It is a control, not a third phenotype in the
story.

Single-mutant fitness alone predicts betaxanthin at $0.151$, so the shared growth component
is real and substantial. Regressing it out of both sides leaves the amino-acid profile at
$0.133$ against $0.176$ without the control on the same genes, so roughly three quarters of
its predictive power survives. Combining the two reaches $0.200$, above either alone.

\notesec{Correction: most of that apparent shrinkage was the gene set, not the control}
Requiring a Costanzo record shrinks the population from 4,432 to 3,708, and the 724 it drops
are not a random quarter. Rather than accept that, every single-deletion fitness source in
the repository was pooled. Kuzmin does \emph{not} help, and the reason is structural: Kuzmin
2018 and 2020 screened the same SGA deletion array Costanzo did, agreeing with it at
$r = 0.947$ to $0.995$ on the overlap, so between them they recover only 56 of the 724.
Baryshnikova 2010 recovers 375, O'Duibhir 2014 a further 167, and the Costanzo NatMX query
arm, which the incumbent control excluded, another 317. The union controls \textbf{4,214 of
the 4,432}, leaving 218.

\begin{table}[htbp]
\centering\small
%% SOURCE: experiments/023-metabolome-betaxanthin-joint/scripts/
%% smf_coverage_fitness_control_sources.py -> results/
\begin{tabular}{lrrrr}
\toprule
gene set & $n$ & 19 AA, no control & fitness alone & 19 AA, fitness controlled \\
\midrule
Costanzo KanMX only (incumbent) & 3{,}708 & $0.1763 \pm 0.0025$ & $0.1508$ & $0.1326 \pm 0.0026$ \\
\textbf{union of all SMF sources} & \textbf{4{,}214} & $\mathbf{0.2884 \pm 0.0014}$ & $0.1332$ & $\mathbf{0.2699 \pm 0.0017}$ \\
\midrule
all shared, no control possible & 4{,}432 & $0.2980 \pm 0.0029$ & -- & -- \\
still uncontrolled after the union & 218 & $0.3150 \pm 0.0312$ & -- & -- \\
\bottomrule
\end{tabular}
\caption[]{The fitness control costs far less than it appeared to. On the incumbent gene set
the control looks like it removes $0.044$ of predictive power ($0.176 \to 0.133$); on the
enlarged set it removes $0.019$ ($0.288 \to 0.270$). Most of the apparent shrinkage was the
population that requiring a Costanzo record selects, not the control itself. Donor screens
are mapped onto the Costanzo scale by least squares on their overlap and standardized within
donor; transfer correlations run $0.951$ (NatMX) and $0.891$ (Baryshnikova) down to $0.511$
(Kuzmin 2020), so the weakest donors carry the most transfer error. The incumbent row
reproduces the committed result to 13 decimal places, which is what makes the two rows
comparable.}
\label{tab:aa-fitness-union}
\end{table}

So the conclusion strengthens rather than weakens: \textbf{$94\,\%$ of the amino-acid
profile's predictive power survives controlling for single-mutant fitness} on the population
where the control can be applied at all.

\notesec{And the deletions no fitness screen reaches carry the extremes} The 724 deletions
with no Costanzo KanMX record have a betaxanthin standard deviation of $0.914$, against
$0.464$ among the 3,708 that do have one, so they are roughly twice as spread out in
production and carry a disproportionate share of the extreme producers. On that same 724 the
amino-acid profile predicts at $0.455$. Even after the union above, 218 remain uncontrolled,
and the profile predicts on those at $0.315$.

\emph{Why they are missing is itself informative, and unverified here.} A gene absent from a
KanMX deletion-collection fitness screen is typically one whose deletion grows too poorly to
score, or that the collection never carried. Either way the missing set is enriched for
strongly-affected strains, which is the plausible reason it carries the extremes. That
mechanism has not been tested.

\tcfigfit{figures/betaxanthin_amino_acid_predictivity.pdf}{(a) Marginal correlation of each
amino acid with betaxanthin over the 4,432 shared deletions; tyrosine highlighted. These
reproduce \file{experiments/019-simb-multimodal/results/pigment_noise_ceiling.json}
exactly, which the generating script asserts rather than assumes. (b) Cross-validated ridge
fits from \cref{tab:aa-predictivity}. (c) Median tyrosine by betaxanthin decile. The
relationship is non-monotone: tyrosine is highest in the lowest-production decile, falls
through the middle, and rises again at the top, which is why a linear correlation reads
near zero. The axis spans $2.65$ to $2.89$ mM, so the swing is about $9\,\%$ of the
median.}{fig:aa-coupling}

\notesec{A mechanism, labeled as unverified} The Cachera background carries
\gene{ARO4}$^{\mathrm{K229L}}$ and \gene{ARO7}$^{\mathrm{G141S}}$, feedback-resistant
alleles that release the shikimate pathway from the control that shapes the pools Mulleder
measured in the plain deletion collection. A plausible reading is that free tyrosine in an
unengineered strain is not the quantity that limits flux in an engineered one. This has not
been tested and no experiment here can test it.

\subsection{Does the metabolome head actually help the betaxanthin head\secstatus{todo}}

%% PLACEHOLDER-023: paired arm analysis, filled from round_leaderboards.csv
The claim under test is narrow and was stated before the round: in Delta grid cell
\file{s04_L6_maskon_lr0.0001}, the joint arm read $+0.0203$ over the control on
\file{val/betaxanthin/pearson_per_feature} at a matched budget of 333 epochs. That was
one paired difference at one seed with no estimate of the noise floor on it. The same grid
read $-0.049$ and $-0.060$ in the two $L = 2$ cells, both run to 1,000 epochs, so the sign
flips with depth.

%% GENERATED by experiments/019-simb-multimodal/scripts/build_retrospective_tables.py
%% SOURCE: results/round_leaderboards.csv, strand betaxanthin_amino_acid_joint.
\begin{table}[htbp]
\centering\small
\begin{tabular}{lrrrrl}
\toprule
grid cell & control & joint & $\Delta$ & aux $r$ & epochs (ctrl/joint) \\
\midrule
\texttt{s00\_L2\_maskon\_lr0.0001} & nan & nan & +nan & -- & nan/nan \\
\texttt{s02\_L2\_maskoff\_lr0.0001} & nan & nan & +nan & -- & nan/nan \\
\texttt{s04\_L6\_maskon\_lr0.0001} & nan & nan & +nan & -- & nan/nan \\
\texttt{s05\_L6\_maskon\_lr0.001} & nan & nan & +nan & -- & nan/nan \\
\texttt{s06\_L6\_maskoff\_lr0.0001} & 0.4261 & nan & +nan & -- & 998/nan \\
\texttt{s07\_L6\_maskoff\_lr0.001} & nan & nan & +nan & -- & nan/nan \\
\texttt{s08\_L4\_maskon\_lr0.0001}$^{\dagger}$ & 0.4149 & 0.4193 & \textbf{+0.0043} & 0.114 & 998/704 \\
\texttt{s09\_L4\_maskon\_lr0.001} & nan & nan & +nan & -- & nan/nan \\
\texttt{s10\_L4\_maskoff\_lr0.0001} & nan & nan & +nan & -- & nan/nan \\
\texttt{s11\_L4\_maskoff\_lr0.001} & nan & nan & +nan & -- & nan/nan \\
\texttt{s01\_L2\_maskon\_lr0.001}$^{\ast}$ & 0.0287 & 0.0413 & \textbf{+0.0125} & 0.026 & 998/998 \\
\texttt{s03\_L2\_maskoff\_lr0.001}$^{\ast}$ & 0.0200 & nan & +nan & -- & 998/nan \\
\midrule
mean, cells where the control learned (9 cells) & & & +nan $\pm$ nan & & 0 of 9 positive \\
mean, all matched-exposure cells (11 cells) & & & +0.0125 $\pm$ nan & & 1 of 11 positive \\
\bottomrule
\end{tabular}
\caption[]{Betaxanthin score with and without a 19-amino-acid auxiliary head, paired within grid cell, on \texttt{val/betaxanthin/pearson\_per\_feature}. $\Delta$ is joint minus control. \emph{aux} is the auxiliary head's own score, carried because an auxiliary head at $r \approx 0$ that still moved the primary metric would mean the gain came from regularization rather than from shared metabolic signal, and those are different findings. $^{\dagger}$ marks a cell whose two arms ran epoch counts differing by more than 50; since the score is a running maximum, such a difference is part effect and part unequal exposure, and those rows are excluded from both means. $^{\ast}$ marks a cell whose control arm scored below 0.10, meaning neither arm learned the task; every one of them is an $\mathrm{lr} = 10^{-3}$ cell, and a difference between two failed runs measures nothing about the auxiliary head.}
\label{tab:bx-aa-paired}
\end{table}


Scored uniformly across the grid, the effect is \textbf{negative and small}, and the
honest summary of it is three sentences rather than three readings.

The mean over the cells whose control arm learned the task at all is negative, the spread
across cells is several times the effect, and every cell carries $n = 1$ per arm except
\texttt{s08}, so there is no within-cell noise floor to measure the difference against.
Depth and the graph mask each move the result by more than the auxiliary head does, which
means any single-cell claim about the head is smaller than the variation across the cells it
could have been made in. And the auxiliary head is itself weak wherever the primary head is
good, scoring well under the $0.209$ the amino-acid strand reaches when it is the only head,
so the joint arm is paying capacity for a head that is underperforming.

\notesec{Read this as one crude attempt, not as a verdict} Adding an unstructured second
head is the least informative version of the question. Nothing in the architecture says the
two phenotypes are connected through a pathway, so the encoder is being asked to discover a
stoichiometric relationship from 19 noisy concentrations and one pigment readout. Given
\cref{tab:aa-incremental}, which shows the information is genuinely there and genuinely
incremental, the finding to carry forward is that \textbf{the route by which a second head's
representation reaches the first head's readout is the thing that failed}, and that a
constraint-based coupling is the version worth trying (\cref{sec:directions}).

The designed replication,
\file{experiments/023-metabolome-betaxanthin-joint/conf/igb_bx_aa.yaml}, adds a third
setting: an auxiliary head restricted to \textbf{tyrosine only}, against a common control,
on a gene set pinned by \file{require_head_targets} so no arm can win by seeing more rows.
\textbf{It has not been run.} Given \cref{tab:aa-predictivity} its tyrosine arm is now
predicted to be the null one and its 19-amino-acid arm the live one, which is a sharper
prediction than the config was written against.

\begin{keybox}
\textbf{The coupling is real, incremental, and not being used.} The free amino-acid profile
of a deletion predicts that deletion's betaxanthin at $r = 0.298$ out of fold, against
$0.064$ for tyrosine, the chromophore precursor, so the axis is the profile and not the
precursor. It is not a growth artifact: pooling every single-deletion fitness screen in the
repository controls 4,214 of the 4,432 shared deletions, and $94\,\%$ of the profile's power
survives that control. It is not redundant with betaxanthin supervision either: $80\,\%$ of
it survives residualizing out a trained betaxanthin model's own held-out predictions, which
correlate with the profile's predictions at only $0.22$. And an unstructured auxiliary head
does not convert any of that into a gain, scoring negative across the grid. \textbf{The
bottleneck is the route from a second head's representation to the first head's readout, not
the availability of the signal.}
\end{keybox}

\notesec{Cannot say} How much of the $0.298$ is attenuated by Mulleder's unmeasurable noise,
so it is a floor rather than an estimate. Whether the redundancy result survives a
\emph{non-linear} control, or a control model stronger than the $0.356$ best among those
that dumped predictions. And whether any architecture converts the coupling into a gain,
since the only mechanism tried is the one least likely to work.

%% notes-tex/019-simb-multimodal/sections/6-carotene.tex
%% [[experiments.019-simb-multimodal.scripts.pigment_noise_ceiling]]
%% SOURCE: ceiling from experiments/019-simb-multimodal/results/pigment_noise_ceiling.json;
%% scores from experiments/019-simb-multimodal/results/round_leaderboards.csv.

\section{Beta-carotene production\secstatus{todo}}
\label{sec:carotene}

\subsection{Why this target is different from the other pigment\secstatus{todo}}

Ozaydin 2013 (doi:10.1016/j.ymben.2012.07.010) screened the deletion collection carrying a
beta-carotene pathway and scored each colony's color \textbf{by eye on an ordinal scale
from $-5$ to $+5$}. The dataset holds 4,474 records, of which 4,344 have a single replicate
and 130 have two or three. That is a different kind of measurement from the Cachera
fluorescence readout, and it needs a different ceiling: a Pearson ceiling on a subjective
ordinal is the wrong object, so rank agreement is used.

Two estimates exist, and they disagree by a factor of seven.

\begin{table}[htbp]
\centering\small
%% SOURCE: experiments/019-simb-multimodal/results/pigment_noise_ceiling.json, beta_carotene
\begin{tabular}{llrr}
\toprule
estimate & basis & $n$ & Spearman \\
\midrule
replicate max against min & 130 replicated rows & 130 & \textbf{0.544} \\
independent re-screen & \file{2ndRoundOfTransformations} sheet & 119 & 0.075 \\
\quad range-restriction corrected & Thorndike case 2 & 119 & 0.105 \\
\bottomrule
\end{tabular}
\caption[]{Reliability of the beta-carotene score. Neither estimate is clean. The
max-against-min comparison is biased in both directions at once: the two are drawn from one
replicate set, which inflates agreement, and they are the widest possible split of that
set, which deflates it. The re-screen is a genuine test-retest but was run on selected top
hits, so its 1st-screen scores sit almost entirely in $3\ldots5$ of a $-5\ldots+5$ scale.}
\label{tab:carotene-ceiling}
\end{table}

\notesec{Why the $0.075$ is not evidence that the score is worthless} The re-screen is the
better \emph{design} of the two, a genuine independent test-retest rather than two draws
from one replicate set. But it was run only on \textbf{selected top hits}, so its first-round
scores sit almost entirely in $3 \ldots 5$ of a $-5 \ldots +5$ scale, with a standard
deviation of $0.89$ against $1.47$ over the full screen (\cref{fig:carotene-ceiling}b).

A correlation computed on a narrowed slice of a population is attenuated by the narrowing
alone, independently of how good the measurement is. The intuition: if every strain
re-screened is already known to be a high producer, the only variation left to correlate is
the variation \emph{among} high producers, which is small and mostly noise. Thorndike's case
2 formula corrects for exactly this, given the restricted and unrestricted standard
deviations, and it takes the $0.075$ to $0.105$. That is still low, but it is a correction
for one restriction, and the sampling is selected on more than range alone.

The reported ceiling is therefore $0.544$, from the replicate comparison, with the caveat
that neither estimate is clean and they disagree by a factor of five even after correction.

\tcfigfit{figures/beta_carotene_noise_ceiling.pdf}{Reliability of the beta-carotene score,
which is rank agreement rather than a Pearson because the target is a hand-scored ordinal.
(a) The primary estimate, comparing the maximum against the minimum score among the 130
strains with more than one replicate, over unrestricted strains: Spearman $0.544$. (b) The
independent re-screen, which is a genuine test-retest but was run only on selected top hits;
the shaded band marks the middle $90\,\%$ of its first-screen scores, which spans a standard
deviation of $0.89$ against $1.47$ over the full screen. That narrowing is the range
restriction responsible for its raw $0.075$, and it is visible here rather than
asserted.}{fig:carotene-ceiling}

\subsection{What was measured\secstatus{todo}}

The strand is 136 runs across three projects. The best validation
\file{spearman_per_feature} is $\mathbf{0.1371}$, run \texttt{5cjpn6zj} in
\file{beta_carotene_v3}, peaking at \textbf{epoch 9 of 49}. The median over all 136 is
$0.0513$. Against the $0.544$ ceiling that is $25\,\%$ realized.

\notesec{Correction, and it is a correction to the method rather than to a judgment call}
Revision 2 reported the best as $0.118$ from run \texttt{o0aadmot} peaking at epoch 799 of
999. That run is only the \emph{third} best in the strand. The leaderboard had fetched
validation curves for a candidate subset of each project, chosen as the top runs by final
logged value and by epoch count, and \texttt{5cjpn6zj} peaks at epoch 9, so it ranks low on
\emph{both} keys at once and was invisible to the rule. \textbf{This is a structural blind
spot, not bad luck}: a run that peaks very early and then decays is exactly what the
candidate rule cannot see, and it falsifies the previous assumption that incomplete run
coverage risked only medians and never maxima (\cref{sec:coverage}).

\begin{keybox}
\textbf{Beta-carotene is limited by its measurement, not by the model.} A hand-scored
ordinal with one replicate per strain has a reliability no architecture recovers, and the
$0.1371$ best \file{spearman_per_feature} sits at $25\,\%$ of a $0.544$ ceiling that is
itself uncertain by a factor of five. The strand carries real signal and is worth keeping
for the coverage argument in \cref{sec:betaxanthin}, where the claim is about which genes a
flux model can represent and does not depend on the absolute score.
\end{keybox}

\notesec{Correction to revision 2 on convergence, which the coverage fix reverses} Revision 2
said the best run ``peaks late (epoch 791 of 999), so it shares the expression strand's
under-training rather than the metabolite strands' early overfitting''. With the true best
run in hand, the opposite is the case: it peaks at \textbf{epoch 9 of 49} and decays. So
beta-carotene groups with the \emph{metabolite} strands, which peak early and then overfit,
not with expression. That is also what makes it a low-value target for more compute: extra
epochs are not what it is short of.

\notesec{Cannot say} Which of the two ceiling estimates is right, and therefore whether
$0.1371$ is a quarter of achievable or most of it. That gap is resolvable only with a
replicate design the source data does not contain.

\notesec{Recommendation} Keep beta-carotene as a second production target for the coverage
argument in \cref{sec:betaxanthin}, where the claim is about which genes a flux model can
represent and does not depend on the absolute score. Do not spend a sweep on lifting its
number.

%% notes-tex/019-simb-multimodal/sections/7-common.tex
%% SOURCE: the ceiling column is assembled from
%% experiments/019-simb-multimodal/results/{expression_ceiling_replicate,
%% morphology_noise_ceiling,pigment_noise_ceiling}.json; the score column from
%% results/round_leaderboards.csv. Both are written by committed scripts.

\section{What the strands have in common\secstatus{todo}}
\label{sec:common}

\subsection{One table, read the same way\secstatus{todo}}

%% GENERATED by experiments/019-simb-multimodal/scripts/build_retrospective_tables.py
%% SOURCE: results/round_leaderboards.csv + the three ceiling JSONs. Do not hand-edit.
\begin{table}[htbp]
\centering\small
\begin{tabular}{lrrrrl}
\toprule
strand & ceiling & best & of ceiling & epochs (peak/last) & runs \\
\midrule
Expression (Kemmeren, Sameith) & 0.775 & \textbf{0.229} & 30\% & 4007/4105 & 163 \\
Expression, masked-label objective & 0.775 & \textbf{0.241} & 31\% & 9188/9997 & 16 \\
Morphology (CalMorph, Ohya) & 0.611 & \textbf{0.082} & 13\% & 27/65 & 23 (1 flat) \\
Expression and morphology, joint & -- & \textbf{0.051} & -- & 28/70 & 32 \\
Betaxanthin (Cachera) & 0.914 & \textbf{0.372} & 41\% & 254/999 & 44 \\
Betaxanthin with metabolome head & 0.914 & \textbf{0.430} & 47\% & 176/496 & 46 \\
Amino-acid pools (Mulleder) & -- & \textbf{0.211} & -- & 402/999 & 31 \\
Beta-carotene (Ozaydin) & 0.544 & \textbf{0.132} & 24\% & 799/999 & 39 \\
\bottomrule
\end{tabular}
\caption[]{Every strand, scored the same way. \emph{best} is \texttt{roll\_max}, the maximum of a centered five-point rolling mean of the validation curve, an upward-biased order statistic whose bias grows with epochs run, which is why the epoch column travels with it. \emph{epochs (peak/last)} gives where the best run peaked and where it stopped: a peak in the last tenth means the run was cut off while still improving. Ceilings are measured by three different estimators because the three measurement types require them; the amino-acid cell is empty because Mulleder has one replicate per strain, and no ceiling is estimable. \emph{flat} counts runs whose validation metric never left zero.}
\label{tab:strand-summary}
\end{table}


\tcfigfit[0.72]{figures/retrospective_achieved_vs_ceiling.pdf}{Achieved score against
measured ceiling, best run per strand. The dashed outline is the ceiling and the gap is the
headroom; amino acid and the expression-morphology joint arm carry no outline because no
single ceiling applies to them. A bare correlation is not comparable across these rows,
which is the reason the panel exists.}{fig:achieved-vs-ceiling}

\tcfigfit[0.72]{figures/retrospective_peak_position.pdf}{Where the best run's peak sits
inside its own budget. Points near the top were still improving when the run stopped, so
their scores are lower bounds set by compute. Points in the lower band peaked early and
spent the rest of the budget overfitting.}{fig:peak-position}

Three things are visible only when the strands sit together. The ceilings differ by more
than a factor of two, so an unqualified $r$ is uninterpretable across rows. The epoch
budgets differ by more than a factor of a hundred, so the score column is partly a record of
how long each strand was allowed to run. And the two strands closest to their ceilings,
betaxanthin and amino acid, are the two with a scalar or low-dimensional target.

\subsection{Two opposite convergence failures, and neither was diagnosed from a curve\secstatus{todo}}

\begin{itemize}
\item \textbf{Expression never peaks.} At every budget tried, up to 10,000 epochs, the
  validation Pearson peaks in the last tenth of the run. Raising the budget from 1,600 to
  10,000 raised the score from $0.204$ to $0.241$ and moved the peak with it.
\item \textbf{The metabolite strands peak early and then overfit.} Amino-acid runs peak at
  epochs 48 to 134 with validation slope turning negative while training slope stays
  positive. The best betaxanthin run peaks at epoch 254 of 999.
\item \textbf{Morphology was never run long enough to tell which it does.} Longest run,
  121 epochs.
\end{itemize}

Same encoder, same optimizer family, opposite behavior. The difference is in the target: a
6,127-dimensional vector of $\log_2$ ratios against a scalar or a 19-vector. Any single
epoch budget applied across strands is wrong for at least two of them, which is what
produced the truncated expression waves and would produce overfit metabolite runs if the
budget were set from expression.

\subsection{Loss and metric point in different directions on expression alone\secstatus{todo}}

Validation loss bottoms 40 to 250 epochs before validation Pearson peaks in wave 4b, and
$\approx 1{,}200$ epochs before it in the long runs, while the prediction
standard-deviation ratio climbs from $0.001$ to $0.52$. The 010 fitness and interaction task
does not do this: there both losses fall monotonically and the validation-loss minimum is at
the last epoch. The consequence already cost the project every checkpoint the expression
strand saved before best-by-metric checkpointing landed.

\notesec{What this suggests, unverified} A quantile loss on a $\log_2$ ratio can be
minimized by a narrow, well-calibrated predictive distribution that gets the ordering
wrong, and Pearson is invariant to the global rescale that the loss is not. Reporting
\texttt{nmse} alongside, defined so that $\texttt{nmse} = 1$ is exactly ``predict each
gene's mean'', separates right-ordering from right-magnitude. That instrument exists; the
experiment that would test whether a different loss closes the gap does not.

\subsection{Configuration choice moves results more than method choice\secstatus{todo}}

On betaxanthin the ten grid cells span Spearman $0.013$ to $0.158$ and AUC $0.406$ to
$0.695$, against a between-method gap of $0.0014$ on AUC. Four low points are the
$\mathrm{lr} = 10^{-3}$ cells, which collapsed to a constant predictor. On the amino-acid
strand the best run is $0.211$ and the median is $0.0498$.

Any single-number claim about a method is dominated by which cell is quoted, which is why
the cell must be fixed by validation in advance, and why the leaderboard records the median
alongside the maximum. It also means an arm difference smaller than the cell spread is not
measurable without pinning the cell and replicating seeds inside it.

\subsection{Cross-modality overlap is real and small\secstatus{todo}}

Before spending compute on joint proteome and expression training, the linear overlap
between the two was measured on the 1,350 strains they share. Per-gene correlation across
strains has median $r = 0.08$, with $1.7\,\%$ of genes above $0.3$ and $8.9\,\%$ negative.
A ridge map recovers held-out $R^2$ of $0.035$ in one direction and $0.024$ in the other,
above a trivial baseline of $-0.004$ but far from redundancy.
\src{experiments/019-simb-multimodal/results/proteome_expression_eda.json}

\tcfigfit{figures/proteome_expression_per_gene_corr_hist_2026-07-22-17-15-26.pdf}{Per-gene
correlation across strains between the Messner proteome and the Kemmeren expression
compendium, over the 1,832 genes and 1,350 strains they share. A narrow lump barely off
zero. The few high-$r$ genes are dominated by the strain in which that gene is itself
deleted, which drives both its transcript and its protein toward zero.}{fig:prot-expr}

The two modalities were measured in different media, synthetic minimal for the proteome and
synthetic complete for the expression, which attenuates every number above. The overlap is
therefore certified as non-redundant and not certified as cleanly complementary. That
distinction is what bounds how hard a ``the proteome helps expression'' claim can be pushed
until a same-media pair exists.

\subsection{The graph prior has never been checked\secstatus{todo}}

The nine interaction graphs enter the model as a hard additive attention mask, so composing
$L$ layers makes the influence of gene $X$ on gene $Y$ approximately the $L$-step
random-walk reachability of $Y$ from $X$ on that graph. That is a falsifiable assumption
about the data: \emph{deleting $X$ perturbs the genes near $X$ on graph $k$}. If it is
false, pulling attention toward the graph pushes the model toward the wrong prior and no
setting of the regularization weight rescues it.

The probe that would test it is specified and unbuilt. It costs minutes of CPU: for each
deleted gene, the across-reporter rank correlation between graph proximity and response
magnitude, reported per graph as an AUC against degree-preserving rewiring and a matched
random graph, where $0.5$ means the graph says nothing.
\src{notes/experiments.019-simb-multimodal.graph-prior-probe.md}

Nine graphs consume nine attention heads. If three carry signal, the other six are
constrained toward noise and are worse than free heads. This is the cheapest unrun
experiment in the project.

\subsection{The perturbation axis is the binding resource\secstatus{todo}}

\Cref{tab:label-accounting} makes the point for expression: 9.2 million scalar labels
across 1,484 perturbations, each gene deleted exactly once. Morphology is the same shape,
1.3 million labels across 4,718 perturbations, each once. The trigenic interaction data is
the opposite shape, 91 thousand labels across 91 thousand records in which a given gene
recurs with many partners.

The strands that work in this project are the ones where the supervision constrains a gene
repeatedly. \textbf{Hypothesis, untested:} the reason multi-task training is attractive here
is not that the labels are informative about each other, but that a second phenotype gives
a second, differently-shaped constraint on the same perturbation. That predicts the gain
should be larger where the two phenotypes share strains and smaller where they merely share
genes, which is a testable statement and has not been tested.

%% notes-tex/019-simb-multimodal/sections/8-directions.tex
%% SOURCE: hand-authored. This section proposes work; it reports no measurement that is not
%% cited to an earlier section.

\section{What to run next\secstatus{todo}}
\label{sec:directions}

Ordered by cost against what each would settle. The first four are cheap enough that not
running them is the expensive choice.

\subsection{Tier 0, cheap and decisive\secstatus{todo}}

\begin{enumerate}
\item \textbf{The graph-prior probe} (\cref{sec:common}). Minutes of CPU, no GPU, no
  training. It asks whether deleting gene $X$ actually perturbs $X$'s neighbors on each of
  the nine graphs, which is the assumption the entire graph channel rests on and which has
  never been checked. Nine attention heads are currently spent enforcing it. If three
  graphs carry signal, regularize those three and free the other six.

\item \textbf{Morphology at full scale.} Drop \file{require_modalities} for the
  morphology-only arm and train on all 4,718 Ohya deletions instead of the 1,440 that also
  carry expression. This is a config change, and it is the only way the $0.0824$ becomes a
  number about morphology rather than about a quarter of morphology. Run it long enough to
  see a peak, which \cref{sec:common} says cannot be assumed from the expression budget.

\item \textbf{The scalar morphology warm-up}, on \file{A113} (ceiling $0.873$, robust CV
  $2.37$) or \file{C124_C} (Medium bud ratio, $0.698$ and $0.886$).
  Cheap, and it calibrates the epoch budget for the 278-feature target. Not colony size:
  colony size is fitness, which is already the project's strongest strand, and the CalMorph
  size features are the ones deletions barely move.

\item \textbf{Apply the post-hoc rescale to the existing best checkpoint.} Multiplying
  predictions by $r/s$ changes no correlation and moves \file{nmse} from $1.010$ to
  $0.944$, taking the model from worse-than-the-mean to better-than-the-mean in squared
  error (\cref{sec:objective}). An evening's work, no training, and it removes an objection
  the figure would otherwise have to answer.

\item \textbf{Bracket the expression peak, by resuming rather than restarting.} No run has
  observed a maximum at any budget up to 10,000 epochs, and a fresh 10,000-epoch run costs
  3.8 days. Resuming the existing best checkpoint is the cheap form of the same experiment,
  and the epoch at which the curve turns is what sets every arm budget afterward.
  \cref{sec:campaign} proposes $E = 4{,}000$ in the meantime, justified by an arm that has
  been observed to peak at 3,921.
\end{enumerate}

\subsection{Tier 1, the two questions worth GPU\secstatus{todo}}

Both are laid out with their epoch budgets, seed counts and slot costs in
\cref{sec:campaign}, which is the operational version of this subsection. In summary:

\begin{itemize}
\item \textbf{Settle the objective first} (\cref{tab:phase-a}). The quantile loss reaches
  its minimum at epoch 463 and rises for the next 9,500 while Pearson climbs, and the model
  never gets \file{nmse} below 1. A mechanism arm evaluated under an objective that is
  optimizing something else cannot separate ``the mechanism does not help'' from ``the
  objective does not reward it''.
\item \textbf{Then the pair term} (\cref{tab:phase-b}). The reference operator provably
  cannot express a dependence on the pair (deleted gene, reporter gene), and masked
  unmasking does not supply one at the point the score is taken. A rank-64 pair term
  already reaches $0.2274$ in 2.0 days against the masked objective's $0.2362$ in 3.8.
\item \textbf{Pair morphology with fitness for the first full-scale joint test}, since
  fitness shares 4,220 strains with morphology against expression's 1,440. Then run the
  expression pairing on the 1,440 as the mechanism test it was designed to be.
\end{itemize}

Substituting tyrosine for the 19 amino acids, as \file{igb_bx_aa.yaml} proposes, is now
predicted to be the \emph{null} arm rather than the live one, since tyrosine alone predicts
betaxanthin at $0.064$ against the profile's $0.298$. Keep it as a negative control, not as
the headline contrast.

\subsection{Tier 2, the structural options\secstatus{todo}}

\notesec{Constraint-based coupling is the mechanism Figure 6 needs} The 023 result is that
the coupling between amino-acid pools and betaxanthin is present in the data at $r = 0.298$
and that a plain auxiliary head does not use it, costing $-0.0265 \pm 0.0159$
(\cref{tab:bx-aa-paired}). That is a structural gap, and an unstructured second head is the
wrong instrument to close it: it asks the encoder to discover a stoichiometric relationship
from 19 noisy concentrations and one pigment readout, with nothing in the architecture
saying the two are connected through a pathway.

A constraint-based layer says it. Penalizing predictions that violate a stoichiometric or
flux balance over yeast-GEM makes the tyrosine-to-betalamic-acid route a property of the
model rather than something it must infer, and it turns the metabolome head from an
auxiliary target into a set of intermediates the constraint actually references.
\textbf{This is the Figure 6 mechanism, and it is what makes the amino-acid panel a step in
an argument instead of a negative result.}

Two facts size it rather than block it. yeast-GEM covers $19.4\,\%$ of the deletions in
the screen and $27\,\%$ of its high producers (\cref{tab:gem-coverage}), so the constraint
is silent on most of the genome and must be a term that switches off outside the model
rather than a hard requirement. And betaxanthin is already the strand closest to its
ceiling, so the headroom the constraint is competing for is the smallest of any strand:
$0.430$ against $0.914$. Expect it to earn its place by making the mechanism legible and by
extending to targets with no screen at all, rather than by moving betaxanthin's number a
lot.

\notesec{Not an expression rescue} The same layer is attractive for expression on the
argument that metabolism constrains transcription. \cref{tab:label-accounting} is the
reason to be cautious: expression's binding constraint is 1,484 perturbations each seen
once, and a metabolic prior does not add perturbations. Gate it on the same evidence
standard the amino-acid coupling met, a measurable relationship in the data first.

\notesec{Reifying the gene node into mRNA and protein} Not yet. The measured case for it is
thin: per-gene mRNA-to-protein correlation across strains has a median of $0.08$, a ridge
map between the two recovers $2$ to $3.5\,\%$ of held-out variance, and the two datasets
were measured in different media, so the decoupling that would justify separate nodes is
partly a condition artifact. Splitting the node multiplies parameters on an axis the data
does not yet resolve. The case changes when a same-media proteome and expression pair
exists, or when perturbations become fine-grained enough that the central-dogma collapse is
what breaks.

\notesec{Masked-label imputation as its own deliverable} The masked-conditioning oracle
reaches cross-study $0.4838$ at $m = 1000$ against the model's $0.238$, and $97.5$ to
$100.6\,\%$ of that gain survives removing a genotype predictor. It is a real capability
and a defensible product surface, and it is not a genotype-to-expression result. Deciding
whether to ship it as a separate thing is a decision, not an experiment.

\subsection{What to do about Figure 3 and Figure 6\secstatus{todo}}

\notesec{Figure 3} The manuscript currently claims $r = 0.543$ for expression and $0.619$
for morphology inside a paragraph marked \texttt{[FILLER -- R3.]}. The measured values are
$0.238$ and $0.082$. Neither placeholder should survive into a draft anyone reads. The
honest version of the panel is the ceiling-relative one: expression at $31\,\%$ of a
measured $0.775$, morphology at $13\,\%$ of a measured $0.611$ \emph{on a quarter of its
data}, with the fitness panel carrying the strong result. \textbf{A $0.24$ per-gene Pearson
is not a Figure 3 headline.}

Two of the three cheapest routes to a better one are not modeling work at all. Morphology
has never seen more than 1,161 of its 4,718 strains in any of 397 runs, so the
full-scale run is a query change and is the largest single change available to the figure.
And the graph-prior probe decides whether nine attention heads are enforcing a premise
nobody has checked. Both land before any arm ladder is worth launching;
\cref{sec:campaign} sizes what follows.

\notesec{Figure 6} The betaxanthin case is the strongest thing in this document and it does
not depend on beating anyone on accuracy. The argument is coverage: yeast-GEM reaches
$19\,\%$ of the screen and misses $73\,\%$ of its high producers, and CGT finds high
producers in the part a flux method cannot represent at $4.4\times$ over base rate.

The last panel is where the figure either becomes an argument or stays a list, and getting
its \emph{structure} right is the open work. The sequence that carries it: setup diagram;
betaxanthin performance against its $0.914$ ceiling; the coverage gap; the capability gap;
the Flux Cone Learning comparison, framed as a tie on the genes both can score; and then
the structural panel. That last one is the constraint-based layer above: the amino-acid
coupling is real in the data and unused by a plain auxiliary head, and a stoichiometric
constraint over yeast-GEM is what would close it. Shown that way the amino-acid material is
the \emph{motivation} for the mechanism rather than a failed joint-training result, which
is the only framing in which a negative belongs in a figure.

Keep the option text in the draw.io boxes. Several of those panels name alternatives that
have not been run, and a composition that still reads when one of them does not arrive is
worth more than a tidier one that does not.

%% notes-tex/019-simb-multimodal/sections/9-campaign.tex
%% SOURCE: cost figures from experiments/019-simb-multimodal/results/
%% expression_objective_diagnosis.json (measured _runtime and epoch counts) and
%% round_leaderboards.csv. The design is proposed, not measured, and says so.

\section{The expression and morphology campaign\secstatus{todo}}
\label{sec:campaign}

The next campaign is expression plus morphology, for Figure 3. What follows is sized in
GPU-days rather than in arms, because the measured cost per run is what makes most possible
designs unaffordable.

\subsection{The budget arithmetic, measured\secstatus{todo}}

\begin{table}[htbp]
\centering\small
%% SOURCE: results/expression_objective_diagnosis.json (_runtime, epoch)
\begin{tabular}{lrrrr}
\toprule
run & epochs & s/epoch & wall clock & peak at \\
\midrule
v9 \file{M_fine} (masked)  & 9{,}999 & 32.9 & \textbf{91.4 h = 3.81 d} & epoch 9{,}674 ($97\,\%$) \\
v8 \file{V_basis64}        & 4{,}106 & 42.0 & 47.9 h = 2.00 d & epoch 3{,}921 ($95\,\%$) \\
\bottomrule
\end{tabular}
\caption[]{One expression run is a multi-day object. Both runs peaked in their last few
percent, so neither bracketed its own maximum.}
\label{tab:campaign-cost}
\end{table}

At $\approx 35$ s/epoch, one GPU delivers $\approx 2{,}500$ epochs per day. So with $G$
GPUs for $D$ days the campaign has

$$\text{arm-slots} \;=\; \frac{G \cdot D \cdot 2{,}500}{E}$$

runs, where $E$ is the per-arm epoch budget. \textbf{Every design decision below is a
consequence of that fraction}, so it should be evaluated with the real $G$ and $D$ before
anything is launched.

\subsection{Set the epoch budget at 4,000, and the reason is measurable\secstatus{todo}}

The previous round's failure was scoring arms at 300 epochs against a curve that had not
turned. The opposite failure is affordable only once: 10,000 epochs is 3.8 days per arm,
which at four GPUs for four days buys four runs and no replicates.

$E = 4{,}000$ is the defensible middle, and not by taste. The v8 run \textbf{peaked at
epoch 3,921} of 4,106, so a pair-rank arm has been observed to reach its maximum inside
that budget. It costs $\approx 1.6$ days per run, so four GPUs for four days buys
$\approx 10$ slots.

\notesec{What $E = 4{,}000$ gives up, stated plainly} The v9 masked run was still climbing
at 9,674 and reached $0.2362$; truncating it at 4,000 would have scored it near $0.21$.
So a 4,000-epoch comparison is fair between arms and \textbf{understates} the absolute
number by roughly $0.02$ to $0.03$. That is the right trade for an arm comparison and the
wrong one for a headline figure, which is why the winning arm gets one long run afterward.

\subsection{Seeds: two per arm, and the noise floor is already known\secstatus{todo}}

Expression's across-initialization spread was measured at $\approx 0.006$ ($n = 8$, fixed
split) in wave 5, an order of magnitude below the betaxanthin strand's $0.030$
(\cref{tab:replicate-noise}), because the split is pinned so \file{seed} varies weights
only. At $\sigma = 0.006$, two seeds give a paired standard error of $\approx 0.004$, so a
$0.015$ effect is nearly $4\sigma$. \textbf{Two seeds per arm is enough}; three would buy
precision the mechanism differences do not need.

\subsection{Phase A, the objective, and why it goes first\secstatus{todo}}

\cref{sec:objective} found that the quantile loss reaches its minimum at epoch 463 and
rises for the next 9,500 while Pearson climbs, and that the model never gets \file{nmse}
below 1. A mechanism arm evaluated under an objective that is optimizing something else
cannot separate ``the mechanism does not help'' from ``the objective does not reward it'',
so the objective is settled first.

\begin{table}[htbp]
\centering\small
\begin{tabular}{llp{62mm}}
\toprule
arm & objective & what a win means \\
\midrule
\texttt{O\_ref}   & quantile (current) & reference \\
\texttt{O\_zmse}  & MSE on per-gene $z$-scored targets & ``predict the mean'' scores 0 rather than winning; already implemented as \file{standardize_per_feature_target} \\
\texttt{O\_corr}  & correlation objective & optimizing the quantity actually scored \\
\bottomrule
\end{tabular}
\caption[]{Phase A. Three arms, two seeds, $E = 4{,}000$: six slots, $\approx 10$ GPU-days.}
\label{tab:phase-a}
\end{table}

\notesec{Report \file{nmse} and the SD ratio on every arm, not just Pearson} The
calibration identity in \cref{sec:objective} means an arm can raise Pearson while leaving
\file{nmse} above 1, and that is a different result from one that fixes both. Both numbers
are already logged.

\notesec{Free, and worth doing on day one} Apply the post-hoc rescale by $r/s$ to the
existing best checkpoint. It changes no correlation and moves \file{nmse} from $1.010$ to
$0.944$. It is an evening's work and it removes the ``loses to the mean'' objection from
the figure regardless of what the campaign finds.

\subsection{Phase B, the pair term, under the winning objective\secstatus{todo}}

The reference operator provably cannot express any dependence on the pair (deleted gene,
reporter gene), and masked unmasking does not supply one at $k = 0$. The ladder is the
mechanism axis.

\begin{table}[htbp]
\centering\small
\begin{tabular}{llp{54mm}}
\toprule
arm & pair rank & note \\
\midrule
\texttt{P\_ref}      & 0  & additive reference \\
\texttt{P\_basis64}  & 64 & the current best-per-compute: $0.2274$ in 2.0 days \\
\texttt{P\_film}     & 90 & diagonal multiplicative, identity at init \\
\texttt{P\_graph}    & -- & masked graph attention applied \textbf{after} the perturbation operator, so the pair term is a function of network distance. \textbf{Not built.} \\
\bottomrule
\end{tabular}
\caption[]{Phase B. Four arms, two seeds: eight slots, $\approx 13$ GPU-days. \texttt{P\_graph}
requires new code and is gated on the graph-prior probe below.}
\label{tab:phase-b}
\end{table}

\notesec{Gate \texttt{P\_graph} on the probe, and run the probe this week} The probe
(\cref{sec:common}) costs minutes of CPU and asks whether deleting gene $X$ actually
perturbs $X$'s neighbors on each of the nine graphs. If the answer is no, \texttt{P\_graph}
is building a mechanism on a false premise and the nine attention heads currently spent
enforcing that premise should be freed instead. \textbf{Nothing in the campaign is more
decision-relevant per CPU-minute}, and it has been specified and unbuilt since 2026.07.26.

\subsection{Morphology runs alongside, and it is a different problem\secstatus{todo}}

Morphology is not a smaller copy of expression. Its target is 278-dimensional rather than
6,127, no morphology run has ever passed 121 epochs, and every one of the 397 runs across
eight projects trained on 1,161 strains (\cref{tab:morph-ntrain}) out of a 4,718-strain
screen.

\begin{enumerate}
\item \textbf{Drop \file{require_modalities} for the morphology-only arm} and rebuild the
  dataset on all 4,718 Ohya deletions. This is a query and config change, and it is the only
  way the $0.0824$ becomes a statement about morphology rather than about a quarter of it.
  \textbf{Do this first}; it is the largest single change available to the figure.
\item \textbf{Measure the per-epoch cost on the first run.} It is not known: the only
  figure on record is a tiny-model smoke at $\approx 47$ s/epoch against expression's
  $\approx 15$ in the same smoke, which would make morphology the more expensive of the two
  despite the smaller target. Until it is measured, the slot arithmetic above does not
  apply to morphology.
\item \textbf{Run the scalar warm-up in parallel}, on \file{A113} (actin\_n\_ratio, ceiling
  $0.873$, robust CV $2.37$) or \file{C124_C}. It calibrates the epoch budget for the
  278-feature target for a fraction of the cost.
\end{enumerate}

\subsection{The joint question, which is the one Figure 3 is about\secstatus{todo}}

Does expression help morphology. The control already exists and is correct:
\file{delta_joint_expr_morph_000.yaml} pins all three conditions to the 1,440 genotypes
carrying both phenotypes, so \texttt{expr\_morph} minus \texttt{morph} isolates the
auxiliary head from a data-quantity difference. It has never been run past 87 epochs.

\notesec{Run it at two scales, and do not conflate them} On the 1,440 shared strains it is
the mechanism test it was designed to be. Paired with \textbf{fitness} instead of
expression it keeps 4,220 of the 4,718 strains, three times the instance count, which is
where the statistical power is. The two answer different questions and both are worth a
slot; the fitness pairing is the one likely to produce a detectable effect.

\notesec{The prior from the metabolite strands is not encouraging, and should be said out
loud} The one auxiliary-head experiment run to completion anywhere in this project is the
metabolome head on betaxanthin, and it reads $-0.0265 \pm 0.0159$
(\cref{tab:bx-aa-paired}). A second phenotype sharing an encoder has not yet been shown to
help any first phenotype here. The expression-morphology pairing has a better prior than
that one, because the two phenotypes share strains rather than merely genes, but the
campaign should be designed to \textbf{measure} the effect rather than to confirm it, which
means the control arm matters more than the joint arm.

\subsection{What the campaign cannot fix\secstatus{todo}}

Expression sits at $0.238$ against a ceiling of $0.775$, and the perturbation axis has 1,484
points with each gene deleted exactly once (\cref{tab:label-accounting}). No objective and
no pair term changes that. If Phase A and Phase B both come back inside the noise floor,
the honest reading is that the binding constraint is the data rather than the model, and
the figure's expression panel should be built around the ceiling-relative statement and the
imputation capability rather than around a genotype-to-expression number.


\end{document}
